% concatenated from BenMPS26.tex
\documentclass[11pt]{article}
% Proper hyphenation.
\usepackage[american]{babel}

% Graphical packages.
\usepackage{graphicx}
\usepackage{comment}
\usepackage{color}
\usepackage{algorithm}
\usepackage{algorithmic}
\usepackage[left=2cm,right=2cm]{geometry}
% Math packages.
\usepackage{amssymb,amsfonts}
\usepackage{amsthm,amsmath}
\usepackage{mathtools}
\usepackage[hidelinks]{hyperref}
%\usepackage{showkeys}
\newcommand {\mat}[1]{\left[\begin{array}{#1}}
\newcommand {\rix}{\end{array}\right]}

\newtheorem{theorem}{Theorem}
 \newtheorem{corollary}[theorem]{Corollary}
\newtheorem{lemma}[theorem]{Lemma}
 \newtheorem{proposition}[theorem]{Proposition}
 \newtheorem{definition}[theorem]{Definition}
\newtheorem{remark}[theorem]{Remark}
\newtheorem{example}[theorem]{Example}
%\newtheorem{algorithm}[theorem]{Algorithm}
\newcommand{\diag}     {\mathop{\rm diag}\nolimits}
\newcommand{\rank}     {\mathop{\rm rank}\nolimits}
\newcommand{\R}{\mathbb R}
\newcommand{\C}{\mathbb C}
%\newcommand{\ker}     {\mathop{\rm ker}\nolimits}
%\newcommand {\proof}{\par\noindent{\it Proof}. \ignorespaces}
\def\eproof{\space
        {\ \vbox{\hrule\hbox{\vrule height1.3ex\hskip0.8ex\vrule}\hrule}}
        \bigskip}
\newcommand{\nnrm}[1]{{\left\vert\kern-0.25ex\left\vert\kern-0.25ex\left\vert #1
		\right\vert\kern-0.25ex\right\vert\kern-0.25ex\right\vert}}
%

\title{Characterization of stability radii for robustly asymptotically  stable dissipative Hamiltonian differential-algebraic systems}


\author{Peter Benner, Volker Mehrmann, Anshul Prajapati, and Punit Sharma}

\begin{document}
\maketitle

\begin{abstract}
    We study linear time-invariant dissipative Hamiltonian differential-algebraic systems.  %arising in a variety of applications.
    We characterize when the systems are robustly asymptotically stable and derive exact conditions and bounds when this property is lost under structure-preserving perturbations.
	\end{abstract}

Keywords: {dissipative Hamiltonian differential-algebraic equation,  structured stability radius, distance to instability, distance to singularity, distance to nearest high index problem}

Msc Classification: {93D20, 93D09, 65F15}


%{\color{red} Possibly Journal of dynamics and differential equations}
\section{Introduction}%
\label{sec:intro}

In this paper, we analyze the robust asymptotic stability under structure-preserving perturbations for the class of \emph{linear time-invariant dissipative Hamiltonian differential-algebraic equation (dHDAE)} systems of the form
%
\begin{equation}\label{dhdae}
    E\dot{x}=(J-R)Qx ,
\end{equation}
%
with coefficient matrices $E, J, R, Q\in \mathbb C^{n,n}$ (the  set of $n\times n$ complex matrices), where the coefficients satisfy
%
\begin{equation}\label{dhprop}
  E^*Q=Q^*E\geq 0,\quad Q^*JQ=-Q^*J^*Q,\quad Q^*RQ= Q^*R^*Q\geq 0.
\end{equation}
%
Here $Q^*$ denotes the conjugate transpose of a matrix $Q$ and $R\, (\geq)> 0$ denotes that $R$ is positive (semi-) definite.
The function $\mathcal H: x\mapsto \frac 12 x^*E^*Qx$, which describes the energy stored in the system, is called the \emph{Hamiltonian} of the system.

Such dHDAE systems arise, e.g. in  structural mechanics and fluid dynamics; see \cite{GudLMS22} for a large list of applications. To illustrate the need for a detailed structured perturbation analysis, we present the following three examples.
%
\begin{example}\label{ex:MDK}{\rm
The first-order representation of an unforced  linear damped mechanical system is given by
%
\begin{equation} \label{eq_brake_dae}
\begin{bmatrix}M &0 \\0 & K \end{bmatrix} \begin{bmatrix} \dot q \\  \dot  p\end{bmatrix} =   \begin{bmatrix}-D  & -K \\ K &0  \end{bmatrix}
\begin{bmatrix} q  \\  p \end{bmatrix},
\end{equation}
%
where $M,D,K\in\mathbb R^{n,n}$ are Hermitian mass, damping, and stiffness matrices with $M,D\geq 0,K>0$; see,
e.g.,~\cite[Chapter 1]{Ves11}. Here, typically $M>0$, but when small and large masses are occurring in the system then $M$ is close to being singular. The matrices $M,K$ are typically finite element matrices and are subject to perturbations such as modeling and discretization errors, while $D$ usually has a lot of uncertainty, since damping and friction is very hard to model.
}
\end{example}

\begin{example}\label{ex:stokes}
{\rm
The space discretization of the unsteady (in)compressible Stokes or linearized Navier-Stokes equations via finite element or finite difference methods typically leads to dHDAE systems of the form
%
\begin{equation}\label{P-instat-op-general}
\begin{bmatrix} M_v & 0 \\ 0 & M_p \end{bmatrix}
 \begin{bmatrix} \dot{v} \\ \dot{p} \end{bmatrix} =
\begin{bmatrix} A & B\\ -B^* & C\end{bmatrix}
\begin{bmatrix} v\\p\end{bmatrix},
\end{equation}
%
where $M_v=M^*_v>0$ is the velocity mass matrix, $M_p=M^*_p\geq 0$ is the pressure mass matrix which is zero  in the case of an  incompressible fluid and very close to zero if the flow is close to being incompressible,
$B$ is the discretized gradient operator (normalized so that it is of full column rank),  $C=C^*\geq 0$ is an often employed stabilization term in numerical analysis that is typically positive definite but of small norm or zero if no stabilization is used, and  $A$ is a diffusion and dissipation matrix that has a positive semidefinite Hermitian part; see, e.g., \cite{EmmM13}. Here, $v$ and $p$ denote the discretized velocity and pressure, respectively. All coefficient matrices are subject to discretization and modeling errors.}
\end{example}	

\begin{example}\label{ex:poro}{\rm
The discretization of the equations that model the deformation of porous media saturated by an incompressible viscous fluid in first-order formulation as in~\cite[Section~3.4]{AltMU21} leads to a dHDAE of the form
%
\begin{align}
\label{eqn:twoField:opMatrix3}
\begin{bmatrix}Y & 0 & 0\\ 0 & A & 0 \\ 0 & 0 & M\end{bmatrix}
  \begin{bmatrix} \dot{w} \\\dot{u}\\ \dot{p} \end{bmatrix}
  = \begin{bmatrix} 0 & -A & D^* \\ A & 0 & 0 \\ -D & 0 & -K \end{bmatrix}
  \begin{bmatrix} w\\ u\\ p \end{bmatrix} ,
\end{align}
%
where $A,M,Y$ are Hermitian positive definite (with $Y$ being  of very small norm), $K$ is typically Hermitian positive semidefinite, and $D$ is general, non-Hermitian. Here, $u$ represents the discretized displacement field, $w$ the associated discretized velocities, and $p$ the discretized pressure.
A common simplification of the model is to set $Y=0$, then $E={\rm diag}(Y,A,M)=E^*\geq 0$ becomes singular.
All coefficient matrices are subject to discretization and modeling errors.
}
\end{example}

The class of dHDAE systems posesses numerous significant geometric and algebraic properties that are nicely encoded in their representation; see~\cite{AchAM21,DalV98,JacZ12,OrtVME02,MehU23}.
In this paper, we focus on the robust asymptotic stability of dHDAE systems. Before we do so, we first recall the definition of robust asymptotic stability for general unstructured linear DAEs, see \cite{ByeN93,DuLM13}.
%
\begin{definition}\label{def:robasyst}
A general homogeneous linear DAE $E\dot x=Ax$ with constant coefficients
$E,A\in \mathbb C^{n,n}$ is called
\emph{robustly asymptotically stable} if the pair $(E,A)$ is regular, i.e. $\det (\lambda E-A)\not \equiv 0$,  of index at most one, i.e. all blocks in the Kronecker canonical form (see Theorem~\ref{th:kcf} below) associated with the eigenvalue $\infty$ have size at most one,  and has all its finite eigenvalues in the open left half of the complex plane.
%We denote the set of robustly stable matrix pairs $(E,A)$ by $\mathcal R(E,A)$ and its boundary by $\partial \mathcal R(E,A)$.
\end{definition}


The following example, which is  modified from an example in \cite{DuLM13}, provides an illustration of the possible difficulties
in the robust asymptotic stability of DAEs under small perturbations.
%
\begin{example}\label{ex:perind}
{\rm
Consider the DAE (it is actually a   dHDAE)
%
\begin{equation}\label{exp1}
\mat {cc} 1 & 0 \\ 0  & 0 \rix \mat {c} \dot x_1 \\ \dot x_2 \rix =\mat {cc} 0 & 1 \\ -1 & 0\rix \mat {c} x_1 \\x_2 \rix,
\end{equation}
%with $Q=I_2$ and $R=0$
which has only the trivial solution $x_1=x_2=0$.

If we perturb (\ref{exp1}) by a small $\varepsilon$ as
%
\begin{equation}\label{exp1p}
\mat {cc} 1 & 0 \\ 0  & 0 \rix \mat {c} {\dot x}_1 \\ {\dot x}_2 \rix =\mat {cc} 0 & 1 \\ -1 & \varepsilon\rix \mat {c} x_1 \\x_2 \rix,
\end{equation}
%
then solving the second equation of (\ref{exp1p}) for $x_2$
and substituting into the first equation, we obtain
%
\begin{equation}\label{uodeexp1}
{\dot x}_1=(1/\varepsilon)x_1.
\end{equation}
%
Clearly, if $\varepsilon>0$, then the perturbed system (\ref{exp1p})  is not a dHDAE any longer, since the Hermitian part of the right hand side $\mat{cc} 0 & 0 \\ 0 &-\epsilon \rix\not \geq 0$. But if $\varepsilon <0$, then the solution  $x_1$ is asymptotically stable, but it qualitatively differs from the solution of the original system (\ref{exp1}). For an arbitrarily prescribed initial value $x_1(0)\neq 0$, the initial value problem for (\ref{uodeexp1}) has a unique solution. Furthermore, the value of $x_2(0)$ is not required and is uniquely determined by $x_1(0)$.
In fact, this small perturbation has changed the index of the dHDAE (\ref{exp1}).
}
\end{example}
%

If the product in \eqref{dhdae} is multiplied out to create a matrix $A=(J-R)Q$ and the dHDAE structure  is ignored, then the asymptotic stability, regularity, and index of the system $E \dot x =A x$ is no longer evident from the structure of the coefficients. See \cite{DuLM13} for a detailed analysis. Also for general linear DAEs it is  very difficult to analyze  robust asymptotic stability under perturbations, see \cite{ByeHM98,GugL25,GugLM17}.
We will show that for dHDAE systems under structure-preserving perturbations this is possible and we present a characterization of  the nearest dHDAE system that is singular, of high index, or has purely imaginary eigenvalues.

As the presented examples demonstrate, for the analysis of models with uncertain data it is important to know whether a dHDAE system is robustly asymptotically stable. Further applications are presented e.g. in~\cite{GraMQSV16,Mar86,MarL90,MarPR07,RomM09,Sch90}. In this context, to quantify which perturbations and errors can be tolerated, it is important to determine the distance to the nearest system that is not robustly asymptotically stable.


For complex unstructured linear systems $\dot x=Ax$ that are asymptotically stable,
%i.e., all its eigenvalues are in the open left half complex plane, the smallest norm of a perturbation that moves an eigenvalue to the imaginary axis is called the \emph{(complex) stability radius}, since arbitrarily small perturbations can then move an eigenvalue to the right half-plane and thus make the system unstable.
%For real systems, there is also the \emph{real stability radius}, which refers to perturbations that are constrained to be real; see~\cite{HinP90}.
the computation of the \emph{stability radius}, i.e., the smallest perturbation that moves an eigenvalue to the imaginary axis, is a very difficult problem, see e.g. \cite{Bye88} and this is even more  the case for unstructured DAE systems, see \cite{ByeN93,DuLM13}.
The situation is  much better for  dHDAE systems. It is well-known \cite{MehMW18} that the finite eigenvalues of the matrix pencil  $\lambda E-(J-R)Q$ associated with \eqref{dhdae} are in the closed left half complex plane, and the purely imaginary eigenvalues (except possibly the eigenvalue $0$) are semisimple. Furthermore, the index of the pencil is at most two and the singular part (if it exists) is easily characterized.
In general, however, the systems are  not necessarily robustly asymptotically stable.



The computation of the smallest perturbation that destroys robust asymptotic stability has recently been studied for ordinary dissipative Hamiltonian systems  (the case that $E=I$) under structure-preserving perturbations; see~\cite{MehMS17,MehMW21}. In~\cite{MehMS16}, bounds for the smallest perturbations were examined by analyzing structure-preserving perturbations to $J$, $R$, and $Q$ individually, resulting in computable formulas. In~\cite{AliMM20}, stability radii approximations for large-scale port-Hamiltonian systems were derived, focusing solely on perturbations in $R$. Additionally, in~\cite{BagGS21}, perturbations to $J$ and $R$ were considered, and a lower bound was derived for the structured stability radii.
In \cite{GilMS18}, the problem  of finding a smallest perturbation that moves a dHDAE system that is not robustly asymptotically stable to the boundary of the set of  robustly asymptotically stable dHDAE systems has been studied.
For the case of dissipative Hamiltonian ODE systems, i.e.\ $E=I$, the distance to instability has recently  been characterized in \cite{BenMPS25}.
%A DAE system $(E,A)$ with nonsingular $E$ can be reformulated as a standard system $(I,AE^{-1})$. However, the reduction process may be numerically unstable. It is shown in~\cite{GilMS18} that for a given pair $(E,A)$ with nonsingular $E$, the distance to the nearest dHDAE $(E,(J-R)Q)$ and the nearest standard DH system $(I,(J-R)Q)$ are significantly different.


In this paper, we study the opposite direction, i.e., we determine formulas and bounds for the distance to the boundary
of the set of robustly asymptotically  stable  dHDAE systems under structure-preserving perturbations.
The paper is organized as follows.
In Section~\ref{sec:prelim}, we present the notation and some preliminary results. Moreover, we  recall some basic properties of dHDAE systems, and we also introduce the concept of robust asymptotic stability.
In Section~\ref{sec:dist}, we present explicit formulas for the structured distance to singularity, the structured distance to the system with index greater than one, and the structured distance to the nearest system with purely imaginary eigenvalues. These results are derived for structured perturbations in all the coefficients. In Section~\ref{sec:partial}, we also present perturbation results for the special cases in which only some of the coefficients are perturbed. Numerical examples are presented in Section \ref{sec:numerical}. We summarize the results and give an outlook on further research problems in Section~\ref{sec:conclusion}.


\section{Notation and preliminaries}\label{sec:prelim}
In the following, $\|\cdot\|$ denotes the spectral norm of a vector or a matrix%, while $\|\cdot\|_F$ denotes the Frobenius norm of a matrix
. By $\Lambda(A)$, we denote the spectrum of a matrix $A \in \C^{n,n}$. %where $\C^{n,r}$ is the set of complex $n \times r$ matrices, with the special case $\C^n = \C^{n,1}$.
%The sets Herm(n) and SHerm(n), respectively, denote the set of complex Hermitian and skew-Hermitian matrices in $\C^{n,n}$.
We use the notation $A \geq 0$ and $A \leq 0$ if $A \in \C^{n,n}$ is Hermitian and positive or  negative semidefinite, respectively, and $A > 0$ if $A$ is Hermitian positive definite.
%For a matrix $A \in \C^{n,r}$ we denote by $A^{\dagger} \in \C^{r,n}$ the Moore–Penrose inverse of $A$; see e.g., [10].
 For a complex number $z$, $\Im(z)$ and $\Re(z)$, respectively, denote the imaginary and the real part of $z$. We denote by $\sigma_{\text{min}}(A)$  the smallest singular value of a matrix $A$. If $R$ is Hermitian, then $\lambda_{\text{max}}(R)$ and $\lambda_{\text{min}}(R)$ denote its largest and smallest eigenvalue, respectively.
We denote the identity matrix of size $n$ by $I_n$, leaving off the index if the dimension is clear from the context.

We will frequently use generalized Rayleigh quotients $\rho(x) := \frac{x^*H_1x}{x^*H_2x}$, with $x\in\C^n\setminus\{0\}$, for Hermitian positive semidefinite $H_1$ and $H_2\in\C^{n,n}$,
where we define that
%
\begin{equation}\label{eq:defrq}
	\rho(x)= \frac{x^*H_1x}{x^*H_2x} :=0, \quad\text{if both $x^*H_1x = 0$ and $x^*H_2x= 0$.}
\end{equation}
%
By \eqref{eq:defrq}, the function $\rho(x)$ is well-defined and lower semi-continuous, i.e., $\liminf_{y\to x}\rho(y)\geq \rho(x)$, for all $x\in\C^n\setminus\{0\}$, so that we can properly define optimization problems involving such Rayleigh quotients.

\subsection{Spectral properties of general DAEs}\label{sec:gendae}

For DAEs $E\dot x=Ax$  we synonymously speak about the pair of matrices $(E,A)$ and the pencil $\lambda E-A$. Then the structural properties are characterized via the Kronecker canonical form \cite{Gan59a}.
%
\begin{theorem}\label{th:kcf}
Let $E,A\in {\mathbb C}^{n,m}$. Then there exist nonsingular matrices
$S\in {\mathbb C}^{n,n}$ and $T\in {\mathbb C}^{m,m}$ such that
%
\begin{equation}\label{kcf}
S(\lambda E-A)T=\diag({\mathcal L}_{\epsilon_1},\ldots,{\mathcal L}_{\epsilon_p},
{\mathcal L}^\top_{\eta_1},\ldots,{\mathcal L}^\top_{\eta_q},
{\mathcal J}_{\rho_1}^{\lambda_1},\ldots,{\mathcal J}_{\rho_r}^{\lambda_r},{\mathcal N}_{\sigma_1},\ldots,
{\mathcal N}_{\sigma_s}),
\end{equation}
%
where the block entries have the following properties:
%
\begin{enumerate}
\item[\rm (i)]
Every entry ${\mathcal L}_{\epsilon_j}$ is a bidiagonal block of size
${\epsilon_j}\times ({\epsilon_j+1})$, $\epsilon_j\in{\mathbb N}_0$,
of the form
\[
\lambda\left[\begin{array}{cccc}
1&0\\&\ddots&\ddots\\&&1&0
\end{array}\right]-\left[\begin{array}{cccc}
0&1\\&\ddots&\ddots\\&&0&1
\end{array}\right].
\]
\item[\rm (ii)]
Every entry ${\mathcal L}^\top_{\eta_j}$ is a bidiagonal block of size
$({\eta_j+1})\times {\eta_j}$, $\eta_j\in{\mathbb N}_0$,
of the form
\[
\lambda\left[\begin{array}{ccc}
1\\0&\ddots\\&\ddots&1\\&&0
\end{array}\right]-
\left[\begin{array}{ccc}
0\\1&\ddots\\&\ddots&0\\&&1
\end{array}\right].
\]
\item[\rm (iii)]
Every entry ${\mathcal J}_{\rho_j}^{\lambda_j}$ is a Jordan block of size
${\rho_j}\times{\rho_j}$, $\rho_j\in{\mathbb N}$, $\lambda_j\in{\mathbb C}$,
of the form
\[
\lambda\left[\begin{array}{cccc}
1\\&\ddots\\&&\ddots\\&&&1
\end{array}\right]-
\left[\begin{array}{cccc}
\lambda_j&1\\&\ddots&\ddots\\&&\ddots&1\\&&&\lambda_j
\end{array}\right].
\]
\item[\rm (iv)]
Every entry ${\mathcal N}_{\sigma_j}$ is a nilpotent block of size
${\sigma_j}\times {\sigma_j}$, $\sigma_j\in{\mathbb N}$,
of the form
\[
\lambda\left[\begin{array}{cccc}
0&1\\&\ddots&\ddots\\&&\ddots&1\\&&&0
\end{array}\right]-
\left[\begin{array}{cccc}
1\\&\ddots\\&&\ddots\\&&&1
\end{array}\right].
\]
\end{enumerate}
%
The Kronecker canonical form is unique up to permutation of the blocks.%, i.e.,
%the kind, size, and number of the blocks are characteristic for
%the matrix pair~$(E,A)$.
\end{theorem}
%
The sizes of the rectangular blocks, $\eta_j$ and $\epsilon_i$,
are called \emph{left and right minimal indices} of $\lambda E-A$, respectively. The matrix pencil $\lambda E-A$ with $E,A\in\mathbb C^{n,m}$ is called \emph{regular} if $n=m$ and
$\operatorname{det}(\lambda_0 E-A)\neq 0$ for some $\lambda_0 \in \mathbb C$,
otherwise it is called \emph{singular}.
%A pencil is singular if and only if it has blocks of at least one of the types ${\mathcal L}_{\epsilon_j}$ or
%${\mathcal L}^\top_{\eta_j}$ in the Kronecker canonical form.


A value $\lambda_0\in\mathbb C$ is called a (finite) eigenvalue of $\lambda E-A$ if
%
\[
\operatorname{rank}(\lambda_0E-A)<\max_{\alpha\in\mathbb C}
\operatorname{rank}(\alpha E-A).
\]
%.
Furthermore,
$\lambda_0=\infty$ is said to be an eigenvalue of $\lambda E-A$ if zero is an eigenvalue of $\lambda A-E$.
The size of the largest block ${\mathcal N}_{\sigma_j}$ is
called the \emph{index} $\nu$ of the pencil $\lambda E-A$, where, by convention,  $\nu=0$ if $E$ is invertible.

% A real Kronecker form  under real transformation
% matrices $S,T$ also exists, where  blocks ${\mathcal J}_{\rho_j}^{\lambda_j}$ with $\lambda_j\in\mathbb C\setminus\mathbb R$ are in real Jordan canonical form instead, but the other
% blocks have the same structure as in the complex case.

For matrix pencils associated with
dHDAE systems we recall the  special Kronecker structure in the next subsection.
\subsection{Spectral Properties of dHDAEs}\label{sec:dhdaes}
The spectral properties of the \emph{dH pencil} $P(\lambda)=\lambda E-(J-R)Q$ associated with~\eqref{dhdae} have been completely characterized in \cite{MehMW18}, see also \cite{GilMS18,MehMS16} for partial results.
%
\begin{theorem}\label{thm:singindspec}
Let $E,Q\in \mathbb C^{n,m}$ satisfy $E^* Q=Q^* E\geq 0$ and let all left minimal indices of
$\lambda E-Q$ be equal to zero (if there are any). Furthermore, let $R\geq 0$. Then the following
statements hold for the pencil $P(\lambda)=\lambda E-(J-R)Q$.
%
\begin{enumerate}
\item[\rm (i)] If $\lambda_0\in\mathbb C$ is a finite eigenvalue of $P(\lambda)$ then $\operatorname{Re}(\lambda_0)\leq 0$.
\item[\rm (ii)] If $\omega\in\mathbb R\setminus\{0\}$ and $\lambda_0=i\omega$ is a finite eigenvalue of $P(\lambda)$, then
$\lambda_0$ is semisimple. Moreover, if the columns of $V\in\mathbb C^{m,k}$ form a basis of a regular deflating
subspace of $P(\lambda)$ associated with $\lambda_0$, then $RQV=0$.
\item[\rm (iii)] The index of $P(\lambda)$ is at most two.
\item[\rm (iv)] All right minimal indices of $P(\lambda)$ are at most one (if there are any).
\item[\rm (v)] If in addition $\lambda E-Q$ is regular, then all left minimal indices of $P(\lambda)$ are zero
(if there are any).
\end{enumerate}
\end{theorem}

It has been demonstrated in \cite{MehMW18} that the   condition $E^* Q=Q^* E\geq 0$ cannot be weakend in general. Furthermore, if  $\lambda E-Q$ is singular,  then the pencil may have eigenvalues with positive real part.

\subsection{Removing the \texorpdfstring{$Q$}{Q} factor in linear pHDAE systems}
\label{sec:noQ}
%
In many applications where dHDAE system arise one has that $Q = I$. In this case $E=E^*\geq 0 $ and  the Hamiltonian is given by $\mathcal H=\frac 12 x^*Ex$. This representation has many advantages, since the coefficients appear linearly in \eqref{dhdae}, which greatly simplifies the analysis and also the perturbation theory. In the following we recall a result from \cite{MehU23} how the factor $Q$ can be removed.

If $Q$ is invertible, then the state equation can be  multiplied with $Q^*$ from the left without changing the solution set, yielding a system with the same solution set given by
%
\begin{equation}\label{mwithQ}
	\widetilde{E}\dot x=Q^* E\dot x = Q^* (J-R)Q x=
 (\widetilde{J}-\widetilde{R})x,   \end{equation}
%
with
$\widetilde{E} \vcentcolon= Q^* E$,
$\widetilde{J} \vcentcolon=Q^* JQ$, $\widetilde{R} \vcentcolon= Q^* R Q$, which
is again a dHDAE, but now has $\widetilde{Q}=I$.
%and hence $\widetilde{E}=\widetilde{E}^*\geq 0$.

If $Q$ is not of full rank, then the situation is more complex and even in the case $E=I$ the solution can grow unboundedly as the following example from \cite{MehMW18} shows. Consider
%
\begin{displaymath}
    \begin{bmatrix} \dot{x}_1 \\ \dot{x}_2 \end{bmatrix} = JQ\begin{bmatrix}x_1\\x_2\end{bmatrix} = \begin{bmatrix} 0&-1\\ 1&0\end{bmatrix}\begin{bmatrix} 1&0\\ 0&0\end{bmatrix} \begin{bmatrix}x_1\\x_2\end{bmatrix},\quad
    \begin{bmatrix}x_1(0)\\x_2(0)\end{bmatrix} = \begin{bmatrix}x_{1,0}\\x_{2,0}\end{bmatrix}
\end{displaymath}
%
with Hamiltonian $\mathcal H=\tfrac {1}{2} x_1^2$. It has the solution $x_1=x_{1,0}$, $x_2=x_{2,0} +t x_{1,0}$ and thus has linear growth and hence is not stable, since the system matrix $JQ$ is a Jordan block of size $2$ at the eigenvalue $0$.

In view of this discussion, to have robust asymptotic stability, we assume in the following that $Q$ is invertible and in this case we may as well just assume that the dHDAE system has the form
%
\begin{equation}\label{dhdae1}
E\dot x=(J-R) x,
\end{equation}
%
with $E=E^*\geq 0$, $R=R^*\geq 0$, $J=-J^*$. Actually one should use this representation  even in the case of ordinary differential equations, see \cite{MehU23} for a detailed discussion.

\subsection{A staircase form for dHDAEs}\label{sec:staircase}


To check whether a system of the form \eqref{dhdae1} is robustly asymptotically stable
one can use the following staircase form under unitary transformations from \cite{AchAM21}.
%
\begin{lemma}\label{lem:tSF}
Let $E,J,R\in\mathbb C^{n,n}$ satisfy $E=E^*\geq 0$, $R=R^*\geq 0$ and $J=-J^*$.
Then there exists a unitary matrix $P\in\mathbb C^{n,n}$, such that
%
\begin{eqnarray}
\hat E :=P^*\ E\ P
&=:&\begin{bmatrix}
E_{1,1} & E_{2,1}^* & 0 & 0 & 0 \\
E_{2,1} & E_{2,2} & 0 & 0 & 0 \\
0  & 0 & 0 & 0 & 0 \\
0  & 0 & 0 & 0 & 0 \\
0  & 0 & 0 & 0 & 0
\end{bmatrix}, \nonumber\\
%
\hat R :=P^*\ R\ P
&=:&\begin{bmatrix}
R_{1,1} & R_{2,1}^* & R_{3,1}^*& 0 & 0 \\
R_{2,1} & R_{2,2} & R_{3,2}^* &0 & 0  \\
R_{3,1} & R_{3,2} & R_{3,3} &0 & 0 \\
0  & 0 & 0 & 0 & 0 \\
0  & 0 & 0 & 0 & 0
\end{bmatrix}, \label{staircase:EJR} \\
%
\hat J :=P^*\ J\ P
&=:&\begin{bmatrix}
J_{1,1} & -J_{2,1}^* & -J_{3,1}^* & -J_{4,1}^* & 0\\
J_{2,1} & J_{2,2} & -J_{3,2}^*& 0 & 0\\
J_{3,1} & J_{3,2} & J_{3,3}&  0 & 0\\
J_{4,1} & 0& 0 & 0 &0\\
 0 & 0 & 0 & 0 & 0
\end{bmatrix}.
\nonumber
\end{eqnarray}
%
The matrices are partitioned in the same way, with (square) diagonal block matrices of sizes $n_1$, $n_2$, $n_3$, $n_4=n_1$, $n_5$.
If present, then $
\begin{bmatrix} E_{1,1} & E_{2,1}^* \\ E_{2,1} & E_{2,2}  \end{bmatrix}>0
$, and the matrices $J_{4,1}$, $J_{3,3}-R_{3,3}$ with $R_{3,3}\geq 0$ are invertible.
\end{lemma}
%\proof
%The proof is presented in \cite{AchAM21} by a constructive  algorithm that is presented in Appendix A.
%\eproof
%The staircase form \eqref{staircase:EJR} is based on a sequence of consecutive rank decisions. In the context of the numerical integration of general linear DAEs one proceeds as follows. Whenever a rank decision (indicated by small singular values) is performed, one takes the decision that in case of doubt one assumes a smaller rank. We adopt the same strategy also for systems of the form \eqref{dhdae1}.

\begin{corollary}\label{cor} Let $E,J,R\in\mathbb C^{n,n}$ with $E=E^*\geq 0$, $R=R^*\geq 0$ and $J=-J^*$  be in staircase form \eqref{staircase:EJR}.

\begin{itemize}
\item [a)] The pencil $\lambda E-(J-R)$ is regular if and only if $n_5=0$ if and only if  $\ker \mat{c} E\\ J\\ R\rix\neq \{0\}$.
\item [b)] The pencil is regular and of index at most one if and only if $n_1=n_4=n_5=0$.
\item [c)] The finite eigenvalues of $\lambda E-(J-R)$ are the eigenvalues of\\ $\lambda E_{2,2}- (-J_{3,2}^*-R_{3,2}^*)(J_{3,3}-R_{3,3})^{-1} (J_{3,2}-R_{3,2})$.
\end{itemize}
\end{corollary}
In \cite{AchAM21} the following refinement of the staircase form \eqref{staircase:EJR} has been presented.
%
\begin{corollary}\label{cor:almostkcf}
Let $E,J,R\in\mathbb C^{n,n}$ satisfy $E=E^*\geq 0$, $R=R^*\geq 0$ and $J=-J^*$.
Then there exist nonsingular matrices  $L,Z\in\mathbb C^{n,n}$, such that
%
\begin{eqnarray}
\tilde E :=L\ E\ Z
&=:&\begin{bmatrix}
E_{1,1} & 0& 0 & 0 & 0 \\
0&  E_{2,2}& 0 & 0 & 0 \\
0  & 0 & 0 & 0 & 0 \\
0  & 0 & 0 & 0 & 0 \\
0  & 0 & 0 & 0 & 0
\end{bmatrix}, \nonumber\\
%
\tilde R :=L\ R\ Z
&=:&\begin{bmatrix}
0 & 0 & 0& 0 & 0 \\
0 & R_{2,2}& 0 &0 & 0  \\
0 & 0 & I_{n_3} &0 & 0 \\
0  & 0 & 0 & 0 & 0 \\
0  & 0 & 0 & 0 & 0
\end{bmatrix}, \label{kronecker:EJR} \\
%
\tilde J :=L\ J\ Z
&=:&\begin{bmatrix}
0 & 0 & 0 & -I_{n_4} & 0\\
0 & J_{2,2} & 0& 0 & 0\\
0& 0 & 0&  0 & 0\\
I_{n_4} & 0& 0 & 0 &0\\
 0 & 0 & 0 & 0 & 0
\end{bmatrix},
\nonumber
\end{eqnarray}
with $E_{11}>0$ diagonal, $E_{22}>0$, and $R_{22}\geq 0$.
\end{corollary}
%\proof
%The proof is given in \cite{AchAM21} followed by a spectral decomposition of the $(1,1) $ block in $E$ and a transformation to bring the $(4,1)$ block and the $(1,4)$ block in $J$ back to $\pm I$.
%\eproof


\begin{comment}
In contrast to the case of general pencils, where the characterization of singularity  is very difficult,  see \cite{ByeHM98}, for dHDAEs of the form \eqref{dhdae1} this can be easily done, \cite{GugM22,MehMW21}.
%
\begin{lemma}\label{lem:regular}
     Consider a
    pencil $\lambda E-(J-R)$ associated with a dHDAE of the form~\eqref{dhdae1}.
  The pencil is regular if and only if $\ker \mat{c} E\\ J\\ R\rix\neq \{0\}$.
    \end{lemma}
\proof
See \cite{MehMW21}.
\eproof
\section{Robust asymptotic stability distances for dHDAEs}\label{sec:robust}

In this section we study the effect of structure-preserving perturbations on dHDAE systems \eqref{dhdae1} of the form
%
\begin{equation}\label{perdhdae}
    (E+\Delta_E)\dot z=(J+\Delta_J -R-\Delta_R)z,
\end{equation}
%
with  $E+\Delta_E=(E+\Delta_E)^*\geq 0$ and $R+\Delta_R=(R+\Delta_R)^*\geq 0$,
$J+\Delta_J=-(J+\Delta_J)^*$, so that the system is still a dHDAE system.
\end{comment}

Based on the spectral properties of dHDAE systems, in the next subsection we give a characterization of robust asymptotic stability.
\subsection{Robust asymptotic stability}\label{sec:robasystab}
The following result gives a characterization of robust asymptotic stability for dHDAE systems.
%
\begin{theorem}\label{thm:stabledh}
    Consider the  dH
    pencil $\lambda E-(J-R)$ associated with the dHDAE~\eqref{dhdae1}. Let $N(E)$ be a matrix with columns that span the right nullspace of $E$. Then the system is robustly asymptotically stable if and only if the following conditions hold:
    \begin{itemize}
        \item [a)] the pencil is regular;
        \item [b)] the index of the pencil is at most one and every principal submatrix of  $N(E)^*(J-R)N(E)$ is nonsingular.
        \item [c)] the maximal real part $\mu$ of a finite eigenvalue of $\lambda E-(J-R)$ (the spectral abscissa) is negative.
    \end{itemize}
    The boundary of the set of robustly asymptotically stable pencils are those pencils where one of the three conditions is violated.
\end{theorem}
\proof
%The proof follows from the unitary staircase form~\eqref{staircase:EJR}.
If the three conditions a)-c) hold, then in \eqref{staircase:EJR} we have that $n_1=n_4=n_5=0$, $E_{22}=E_{22}^*>0$ and $J_{33}-R_{33}$ is nonsingular. Sufficiently small structured perturbations will not destroy the nonsingularity of these two matrices. So it remains to study the case that the rank of $E$ is increased. Since all principal submatrices of $N(E)^*(J-R)N(E)$ are nonsingular, an increase of the rank of $E$ will again lead to a system of index at most one.

For the converse direction we show that a violation of any of the three conditions yields a pair on  the boundary of the robustly asymptotically stable dH pencils.

a) If the pencil is singular, then it has
a zero row and column in the staircase form. Any arbitrary small positive perturbation in the corresponding
diagonal position of $E$ leads to an eigenvalue $0$ of the pencil, hence the pair is on the boundary of the robustly asymptotically stable dH pencils.

b) If the index is bigger than  one, then in \eqref{kronecker:EJR} we have $n_1=n_4>0$. Looking at the structure of the subpencil consisting  of the first and forth block row and column, we have $n_4$ copies of
Example~\ref{ex:perind} and since $n_4>0$ a similar perturbation  will put the pencil on the boundary of the asymptotically stable systems.

c) If $\mu=0$ then the system already has purely imaginary eigenvalues.
\eproof

\begin{remark} \label{rem:R22posdef}{\rm
A sufficient condition for the second part of b) in Theorem~\ref{thm:stabledh} is that $N(E)^*RN(E)>0$ because then every principal submatrix has this property as well.
}
\end{remark}
By structured perturbations of dHDAEs, finite eigenvalues in the open left half plane may move to the imaginary axis, and the regularity or the index of the system may change. %, see Example~\ref{ex:perind}.
In the next subsection we present distance measures under structured perturbations that characterize when this happens.
%
\subsection{Structured distances}\label{sec:structdist}\

In this subsection, we consider perturbations in the coefficient matrices $E$, $J$ and $R$ of a dHDAE system of the form~\eqref{dhdae1}. These take the form
%
\begin{equation}
    (\tilde E, \tilde J-\tilde R)=( E+\Delta_E, J+\Delta_J - (R+\Delta_R)).
\end{equation}
%
In order to measure these perturbations in $E$, $J$ and $R$, we consider the following norm for a tuple $(\Delta_E,\Delta_J,\Delta_R)\in (\C^{n,n})^3$ given by
%
\begin{equation}
    \nnrm{(\Delta_E,\Delta_J,\Delta_R)}=\sqrt{\|\Delta_E\|^2+\|\Delta_J\|^2+\|\Delta_R\|^2}.
\end{equation}
%
%Note that we could also measure the perturbations in an energy norm weighted by $Q$, i.e. $\|x\|_Q=\sqrt{x^HQx}$.
%if $Q$ is positive definite, or the similarly defined semi-norm if $Q$ is only semidefinite.
%The results is the same if we consider pencils of arising from a \eqref{dhdae} with $Q=I$.
We study the following structured distances.


\begin{definition}\label{def:distances}
	Consider a robustly  asymptotically stable dHDAE system of the form~\eqref{dhdae1}, and the following sets  of structured perturbations to the matrices $E, J$ and $R$,
	%
	\begin{align}
		\mathcal S_d(E,J,R)&:= \{(\Delta_E,\Delta_J,\Delta_R):\Delta_E\leq 0, E+\Delta_E\geq 0, \Delta_J^*=-\Delta_J, \Delta_R\leq 0, R+\Delta_R\geq 0 \}, \label{set:Sd}\\
		\mathcal S_i(E,J,R)&:= \{(\Delta_E,\Delta_J,\Delta_R): \Delta_E^*=\Delta_E,E+\Delta_E\geq 0, \Delta_J^*=-\Delta_J, \Delta_R^*=\Delta_R, R+\Delta_R\geq 0 \}.\label{set:Si}
	\end{align}
	%
	Let $\mathbb S\in \{\mathcal S_d,\mathcal S_i\}$. Then, we consider the  distances:
	%
	\begin{enumerate}
		\item The distance to singularity $d_{sing}^{\mathbb S}(E,J,R)$ under perturbations from  $\mathbb S\in \{\mathcal S_d,\mathcal S_i\}$, is defined by
		%
		\begin{equation}\label{def:dist-sing}
			d_{sing}^{\mathbb S}(E,J,R) := \inf\{\nnrm{(\Delta_E,\Delta_J,\Delta_R)}~:~(E+\Delta_E,J+\Delta_J-R-\Delta_R) \text{ is singular}\}.
		\end{equation}
		%
		\item The distance to higher index $d_{hi}^{\mathbb S}(E,J,R)$ under perturbations from  $\mathbb S\in \{\mathcal S_d,\mathcal S_i\}$, is defined by
		%
		\begin{equation}\label{def:dist-hi}
			d_{hi}^{\mathbb S}(E,J,R) := \inf\{\nnrm{(\Delta_E,\Delta_J,\Delta_R)}~:~(E+\Delta_E,J+\Delta_J-R-\Delta_R) \text{ has index }>1\}.
		\end{equation}
		%
		\item The distance to a system with eigenvalues on the imaginary axis $d_{im}^{\mathbb S}(E,J,R)$ under perturbations from the set $\mathbb S\in \{\mathcal S_d,\mathcal S_i\}$, is defined by
		%
		\begin{equation}\label{def:dist-im}
			d_{im}^{\mathbb S}(E,J,R) := \inf\{\nnrm{(\Delta_E,\Delta_J,\Delta_R)}~:~\Lambda(E+\Delta_E,J+\Delta_J-R-\Delta_R)\cap i\mathbb R \neq \emptyset \}.
		\end{equation}
	\end{enumerate}
\end{definition}

By  Theorem~\ref{thm:stabledh}, the structured distance  of a robustly asymptotically stable dHDAE system to the boundary of this set is then given by
%
\begin{equation*}\label{def:dist-inst}
	d_{inst}^{\mathbb S}(E,J,R):= \min \{d_{sing}^{\mathbb S}(E,J,R), d_{hi}^{\mathbb S}(E,J,R),d_{im}^{\mathbb S}(E,J,R)\}.
\end{equation*}
%
%That is, instability is reached by the smallest structure-preserving perturbation that induces either a singularity of the pencil, an increase in the index, or the appearance of a finite eigenvalue on the imaginary axis.

In addition to full perturbations of the dHDAE system, we will also consider partial perturbations, e.g.
%when computing the above distances. The corresponding perturbation sets and structured distances are defined similarly. Specifically,
the partial perturbation sets $\mathcal S_d(J,R):=\mathcal S_d(0,J,R)$ and $\mathcal S_i(J,R):= \mathcal S_i(0,J,R)$.
%, where only $J$ and $R$ are perturbed. These are obtained as subsets of $\mathcal S_d(E,J,R)$ and $\mathcal S_i(E,J,R)$, respectively, by restricting to perturbations with $\Delta_E=0$.
%That is,
%
%\begin{equation}
%	\mathcal S_d(J,R)=\{(\Delta_J,\Delta_R)~:~(0,\Delta_J,\Delta_R)\in S_d(E,J,R)\}, \quad \mathcal S_i(J,R)=\{(\Delta_J,\Delta_R)~:~(0,\Delta_J,\Delta_R)\in S_i(E,J,R)\}.
%\end{equation}
%
%The corresponding distances to singularity, higher index, and the imaginary axis are then defined in an analogous manner.
%by minimizing over perturbations in $\mathbb S \in \{S_d(J,R), S_i(J,R)\}$.



\subsection{Checking robust asymptotic stability}\label{sec:distalg}

Based on  Theorem~\ref{thm:stabledh},
we have the following procedure to check, via a numerically robust algorithm, whether for a dHDAE of the form \eqref{dhdae1} one of the conditions  for robust asymptotical stability is violated. This procedure also leads  to bounds for the distance to the nearest dHDAE which is on the boundary of the set of robustly asymptotically stable dHDAEs.

%\begin{algorithm}
{\bf \label{alg:dist} Algorithm to check robust asymptotic stability}

Step 1: Check whether $E$ is invertible by computing the %Cholesky facto\-rization
%$E=L_EL_E^*$ or the
spectral decomposition $E=Q_e \Lambda_E Q_e^*$ and if $E$ is singular then determine the right nullspace $N(E)$.\\
The distance to the nearest dHDAE system with infinite eigenvalues then is   $d_{DAE}^{\mathbb S}:=\lambda_{\min}(E)$.

Step 2: Check the regularity of the pencil by computing the singular value decomposition of %
\[
U_s^*\mat{c} E\\ J\\ R\rix V_s=\Sigma_s.
\]
The distance to the nearest singular dissipative Hamiltonian system then is
\[
d_{sing}^{\mathbb S}:=\sigma_{\min}\left (\mat{c} E\\ J\\ R\rix \right ).
\]
This distance can also be determined via the algorithm in \cite{GugM22}.

Step 3: Check the index by computing the singular value decomposition of
\[
U^*_iN(E)^*(J-R) N(E)V_i=\Sigma_i.
\]
The distance to the nearest system of index two is then bounded via $d_{hi}^{\mathbb S}\geq \sigma_{\min}(N(E)^* (J-R) N(E))$, see Theorem~\ref{thm:dist-hiEJR} for an exact characterization.



Step 4. Compute the spectral decomposition
%
\[
U_r^*N(E)^*RN(E)U_r=\Lambda_r.
\]
If $\Lambda_r$ is nonsingular, then  $\Delta_R=-\lambda_{\min} (\Lambda_r) U_rx_r x_r^*U_r^*$ is the smallest perturbation such that $R+\Delta_R$ is singular. Here $x_r$ is a normalized eigenvector to the minimal eigenvalue $\lambda_{\min}(\Lambda_r)$.\\
The distance to the nearest system with singular $R$ coefficient is then $d_r^{\mathbb S}:=\lambda_{\min}(\Lambda_r)$.\\

Step 5. Compute the Schur complement pencil associated with the finite eigenvalues
from the staircase form  in Lemma~\ref{lem:tSF} or the orthogonal condensed form in \cite{MehX15}. Let
%
\[
P_f(\lambda):=\lambda E_f-(J_f-R_f)=\lambda E_{2,2}- (-J_{3,2}^*-R_{3,2}^*)(J_{3,3}-R_{3,3})^{-1} (J_{3,2}-R_{3,2}).
\]
%


Determine the smallest perturbation from the corresponding sets
$\mathcal S_d(E_f,J_f,R_f)$, $\mathcal S_i(E_f,J_f,R_f)$
associated with the reduced pencil $\lambda E_f-(J_f-R_f)$ for which an eigenvalue of $P_f$ reaches the imaginary axis
\[
d_{im}^{\mathbb S}(E_f,J_f,R_f).
\]
%
This procedure directly determines an upper bound for the boundary of the set of robustly asymptotically stable dHDAE systems, since we can definitely destroy the robust asymptotic stability by a perturbation of the size $
\min(d_{im}^{\mathbb S},d_{sing}^{\mathbb S},d_{im}^{\mathbb S})$.
The different bounds are related via  the inequalities
%
\begin{eqnarray*}
&& \quad \max (\sigma_{\min} (E),\sigma_{\min} (N(E)^*(J-R)N(E)),d_r^{\mathbb S})\leq d_{hi},
\\
&&\max (\sigma_{\min} (E),\sigma_{\min} (J),\sigma_{\min} (R))\leq d_{sing}^{\mathbb S} = \sigma_{\min} \left (\mat{c}  E\\ J\\ R\rix \right ).
%\\
%&& \mbox{\rm bound in} \cite{MehMW21}.
\end{eqnarray*}
%
We also have the following sufficient
conditions.
%
\begin{theorem}\label{thm:Rposdef}
  Consider the pencil $\lambda E-(J-R)$ associated with the dHDAE~\eqref{dhdae1} subject to a structured perturbation $(\Delta_E, \Delta_J, \Delta_R)$.
If $R>0$ and $R+\Delta_R>0$, then the perturbed dHDAE system is robustly asymptotically stable.
\end{theorem}
\proof
Since $J+\Delta_J-R-\Delta_R$ is invertible if $R+\Delta_R$ is positive definite, the pencil is robustly regular and of index at most one, since also every principal submatrix of $R+\Delta_R$ is positive definite, see Remark~\ref{rem:R22posdef}.
Let $\hat N$ be a matrix with columns that span the orthogonal complement of $N(E+\Delta_E)$.
Then the subpencil associated with the finite eigenvalues, given by
%
\[
\lambda \tilde E-(\tilde J-\tilde R):=\lambda {\hat N}^*(E+\Delta_E) {\hat N}-{\hat N}^*(J+\Delta_J-R-\Delta_R) {\hat N},
\]
%
has only eigenvalues with negative real part, on the left of the line
\[-\lambda_{\min}(\tilde E^{-1/2})^*\tilde R(\tilde E^{-1/2})i.\]
Hence,  the perturbed pencil is robustly asymptotically stable.
\eproof

Thus, to make a dHDAE system to be not robustly asymptotically stable we need to choose $\Delta_R$ at least such that $R+\Delta_R$ is singular.

Another easy bound in the case of real dHDAE systems is presented in \cite{MehMW21}.
%, where for a matrix $Y\in\mathbb R^{n,n}$
%and a vector $u\in\mathbb R^n$ with $\|u\|=1$,  the matrix
%
%\begin{equation}\label{deltaY}
%\Delta_Y^u=-uu^T Y-Yuu^T+uu^TYuu^T
%\end{equation}
%
%is introduced.
%
%\begin{theorem}\label{thm:singind}
%Consider a  pencil  $\lambda E -(J-R)$ associated with a real dHDAE.
%Then the following statements hold.
%
%\begin{enumerate}
%\item
%The distance to singularity is attained with a perturbation
%$\Delta_E=\Delta_E^u$, $\Delta_J=\Delta_J^u$, and $\Delta_R=\Delta_R^u$ as in~\eqref{deltaY}
%for some $u\in\mathbb R^n$ with $\|u\|=1$. The distance is given by
%
%\[
%d_{sing}^{\mathbb S}
%=\min_{u \in\mathbb R^n, \|u\|=1} \sqrt{2\|{Ju}\|^2+2\big\|(I-uu^T)Eu\big\|^2+(u^T Eu)^2+
%2\big\|(I-uu^T)Ru\big\|^2+(u^T Ru)^2}
%\]
%\end{equation*}
%
%and is bounded via
%
%\begin{equation}\label{lmino}
%\sqrt{\lambda_{\min} (-J^2+R^2+E^2)}\leq d_{sing}^{\mathbb S}\leq  \sqrt{2\cdot\lambda_{\min}(-J^2+R^2+E^2)}.
%\end{equation}
%
%\item The distances $d_{im}^{\mathbb S}$ and  $d_{hi}^{\mathbb S}$ coincide and
%satisfy
    %
%    \[
%d_{im}^{\mathbb S}=d_{hi}^{\mathbb S}
%=\min_{u \in\mathbb R^n, \|u\|=1} \sqrt{2\big\|(I-uu^T)Eu\big\|^2+(u^T Eu)^2+
%2\big\|(I-uu^T)Ru\big\|^2+(u^T Ru)^2}
%\]
%
%and are bounded as
%
%\begin{equation*}
%\sqrt{\lambda_{\min}(E^2+R^2)}\leq d_{hi}^{\mathbb S}=d_{im}^{\mathbb S}
%\leq  \sqrt{2\cdot\lambda_{\min}(E^2+R^2)}.
%\end{equation*}
%\end{enumerate}
%\end{theorem}
We will extend this result to the complex case below.

\subsection{Mapping results}\label{sec:mapping}
In this subsection, we recall some mapping results from the literature in a form that allows a direct application in computing the structured distances in this paper.
The following result from~\cite{MacMT08} gives minimal norm solutions to the Hermitian mapping problem with respect to the spectral norm.

%We end this section with the following mapping result \cite{??}.
%
\begin{theorem}\label{map:herm}
	Let $x \in \C^n\setminus\{0\}$ and $y\in \C^{n}$. Then
	%
	\begin{itemize}
		\item [a)] there exists a Hermitian matrix $H \in \C^{n,n}$ such that $Hx = y$ if and only if $\Im{(x^*y)} = 0$ and we have
		\[
		\min\big\{\|H\|\; : H\in \C^{n,n},\; H^*=H, \; Hx=y\big\} = \frac{\|y\|}{\|x\|},
		\]
		and the minimum is attained by
		%
		\begin{equation}\label{def:hatH}
			\hat H_{(x,y)} := \frac{\|y\|}{\|x\|}\mat{cc}\frac{y}{\|y\|}&\frac{x}{\|x\|} \rix \mat{cc}\frac{y^*x}{\|x\|\|y\|}&1\\1&\frac{x^*y}{\|x\|\|y\|} \rix \mat{cc}\frac{y}{\|y\|}&\frac{x}{\|x\|} \rix^*
		\end{equation}
		%
		if $x$ and $y$ are linearly independent and by $\hat H_{(x,y)}:=\frac{yx^*}{x^*x}$, otherwise;
		\item [b)] there exists a skew-Hermitian matrix $S \in \C^{n,n}$ such that $Sx = y$ if and only if $\Re{(x^*y)} = 0$ and we have
		\[
		\min\big\{\|S\|\; :\; S ^*=-S, \; Sx=y\big\} = \frac{\|y\|}{\|x\|},
		\]
		and the minimum is attained by $\hat S:=-i\hat H_{(x,iy)}$, where $\hat{H}$ is defined in~\eqref{def:hatH}.
	\end{itemize}
\end{theorem}

We close this section with the following result~\cite[Theorem 2.3]{MehMS16} that gives minimal norm solutions to the Hermitian positive semidefinite mapping problem with respect to the spectral norm.
\begin{theorem}\label{map:def}
    Let $x,y \in \C^n \setminus \{0\}$. Then there exists a positive semidefinite Hermitian matrix $H$ such that $Hx=y$ if and only if $x^*y > 0$. If the latter condition is satisfied, then
    \[
    \min\left \{\|H\|:H \in \C^{n,n},H^*=H \geq 0,~Hx=y\right \}=\frac{{\|y\|}^2}{x^*y}
    \]
    and the minimum is attained for the rank one matrix
    $\tilde H=\frac{1}{x^*y}yy^*$.
\end{theorem}

\noindent In this section we have summarized some preliminary results as well as the characterization of robust asymptotical stability. In the following section we  characterize the distances in Definition~\ref{def:distances}.

\section{Distance to the nearest dHDAE that is not robustly asymtotically stable under full perturbations}\label{sec:dist}

In this section we discuss in detail the various  distances in Definition~\ref{def:distances}. We derive explicit characterizations as well as bounds.

\subsection{Distance to nearest dHDAE with purely imaginary eigenvalues}\label{sec:distim}
In this section, we present analytic results that characterize the smallest structured distance to a dHDAE system of the form \eqref{dhdae1} with purely imaginary eigenvalues.

We use the following notation. Let $\Lambda:=\{\omega:\det(i\omega E-J) =0\}$, and for $\omega\in \R\setminus \Lambda$, define $M=M(\omega):=(i\omega E-J)^{-1}$ and the  block matrices
%
\begin{equation}\label{mat:tildeH}
    \tilde H_1:=\mat{cc} 0 & -M^*\\ M & 0 \rix, \quad \tilde H_2:=\mat{cc} 0 & -i\omega M \\ i\omega M^* & 0 \rix,
\end{equation}
%
as well as
%
\begin{equation}\label{mat:G1G2}
    G_1:=\mat{c} I\\i\omega I \rix M^*R^2M \mat{cc} I & -i\omega I \rix + \mat{cc}I & \\ & I \rix, \quad G_2:=\mat{c} I\\i\omega I \rix M^*M \mat{cc}I & -i\omega I \rix.
\end{equation}
%
Note that  it follows from Theorem~\ref{thm:Rposdef} that $G_1$ is positive definite, and hence has a unique positive definite square root $G_1^{1/2}$. Thus, we can define the following matrices
%
\begin{equation}\label{mat:GH1H2}
    G:=G_1^{-1/2}G_2G_1^{-1/2}, \quad H_1:=G_1^{-1/2}\tilde H_1G_1^{-1/2}, \text{ and } H_2:= G_1^{-1/2}\tilde H_2G_1^{-1/2}.
\end{equation}
%
For our analysis, we will make use of the following Lemma.
%
\begin{lemma}\label{lem:tfae-map}
    Consider a robustly asymptotically stable dHDAE system of the form~\eqref{dhdae1}, and let $\mathbb S\in \{\mathcal S_d,\mathcal S_i\}$, where $\mathcal S_d,\mathcal S_i$ are defined in~\eqref{set:Sd} and~\eqref{set:Si}, respectively. Furthermore, let $\omega\in \R$ be such that $i\omega$ is not an eigenvalue of the pair $(E,J)$ and define $M:=(i\omega E-J)^{-1}$. Then for $(\Delta_E,\Delta_J,\Delta_R)\in\mathbb S$, the following are equivalent.
    %
    \begin{enumerate}
     \item[\rm (1)] ${\rm det}\left(i\omega (E+\Delta_E)-(J+\Delta_J)+(R+\Delta_R)\right)=0$.
        \item[\rm (2)] $\exists~x(\neq 0)\in \C^{n}$ s.t. $(R+\Delta_R)x=0$, $(i\omega (E+\Delta_E)-(J+\Delta_J))x=0.$
        \item[\rm (3)] $\exists~v_E,v_J\in \C^{n}$~ s.t. $v:=-i\omega v_E+v_J\neq 0$, ~$\Delta_R M v = -RM v, \Delta_EM v=v_E$, and $\Delta_JM v=v_J$.
    \end{enumerate}
\end{lemma}
\proof
The equivalence $(1)\iff (2)$ follows immediately from Theorem~\ref{thm:singindspec}.

$(2)\implies(3)$. If $(2)$ holds, then with $v_J=\Delta_J x$ and $v_E=\Delta_Ex$ we have
%
\begin{equation*}
 (i\omega E-J )x=-i\omega \Delta_E x+\Delta_J x=-i\omega v_E+v_J=:v
\end{equation*}
Since $M=(i\omega E-J)^{-1}$ is nonsingular, we have $v\neq0$ and $x=M v$. This implies that $\Delta_R M v = -RM v$, $\Delta_EM v=v_E$, and $\Delta_JM v=v_J$.

$(3)\implies(2)$ Suppose that $(3)$ holds and set $x=Mv$. Then clearly $(R+\Delta_R)x=RMv+\Delta_R Mv=RMv-RMv=0$ and
\[
(i\omega (E+\Delta_E)-(J+\Delta_J))x= (i\omega E-J)Mv + (i\omega \Delta_EMv-\Delta_JMv)=v-v=0.
\]
This implies (2).
\eproof

Using Lemma~\ref{lem:tfae-map} we have the following bounds for the distance to the nearest dHDAE system with purely imaginary eigenvalues.
%
\begin{theorem}\label{thm:dist-imagEJR}
    Consider a robustly asymptotically stable dHDAE system of the form~\eqref{dhdae1}. Then,
    %
    \begin{enumerate}
        \item  for  perturbations from $\mathcal S_d$, the distance to a dH pencil with purely imaginary eigenvalues
        %the imaginary axis $d_{im}^{\mathcal S_d}(E,J,R)$
        is bounded by
        \begin{align}\label{Sd:dist-imagEJR}
            & (d_{im}^{\mathcal S_d}(E,J,R))^2\geq \min\Bigg\{ \inf_{x\in \Omega}\frac{\|Rx\|^4}{(x^*Rx)^2},\nonumber \\ & \inf_{\omega\in \R\setminus \Lambda}  \left\{\inf_{\substack{u\in\C^{2n},\\Lu\neq 0}}\left(\frac{\|RMLu\|^2}{u^*L^*M^*RMLu}\right)^2+\frac{\|[I_n~0] u\|^2}{\|MLu\|^2} + \left(\frac{\|[0~I_n] u\|^2}{u^*L^*M^*[0~I_n] u}\right)^2: \begin{matrix} u^*[I_n~0]^TMLu \in i\R,\\ ~u^*L^*M^*[0~I_n]u>0 \end{matrix} \right\}  \Bigg\},
        \end{align}
        where $\Omega$ is the set of eigenvectors corresponding to the eigenvalues of the pair $(iE,J)$, and $L:=[I~-i\omega I]$,
        \item for  perturbations from $\mathcal S_i$, the distance to a dH pencil with purely imaginary eigenvalues
        %the imaginary axis $d_{im}^{\mathcal S_d}(E,J,R)$
        is bounded by
        \begin{equation}\label{Si:dist-imagEJR}
            (d_{im}^{\mathcal S_i}(E,J,R))^2\geq \min\left\{ \inf_{x\in \Omega}\frac{\|Rx\|^2}{\|x\|^2}, \inf_{\omega\in \R\setminus \Lambda} \left(\min_{t_1,t_2 \in \R} \lambda_{\max}(G+t_1H_1+t_2H_2) \right)^{-1} \right\},
        \end{equation}
        %
        where $\Omega$ is the set of eigenvectors corresponding to the eigenvalues of the pair $(iE,J)$, where $G$, $H_1$, and $H_2$ are defined in~\eqref{mat:GH1H2}. %Furthermore, equality holds in~\eqref{Si:dist-imagEJR} if the optimal perturbations $\hat \Delta_E$ and $\hat \Delta_R$ in $E$ and $R$, satisfy, $E+\hat \Delta_E\geq 0$ and $R+\hat \Delta_R\geq 0$, respectively.
    \end{enumerate}
\end{theorem}
\begin{proof}
    For $\mathbb S\in \{\mathcal S_d,\mathcal S_i\}$,  the distance to a dHDAE with eigenvalues on the imaginary axis is defined by
    %
    \begin{equation*}
        d_{im}^{\mathbb S}(E,J,R) = \inf\{\nnrm{(\Delta_E,\Delta_J,\Delta_R)}~:~(\Delta_E,\Delta_J,\Delta_R)\in \mathbb S,~\Lambda(E+\Delta_E,J+\Delta_J-R-\Delta_R)\cap i\mathbb R \neq \emptyset \}.
    \end{equation*}
    %
    Using Lemma~\ref{lem:tfae-map}, we obtain
    %
    \begin{align*}
        &d_{im}^{\mathbb S}(E,J,R)\nonumber\\
        &= \inf\{\nnrm{(\Delta_E,\Delta_J,\Delta_R)}:(\Delta_E,\Delta_J,\Delta_R)\in \mathbb S, (R+\Delta_R)x=0 \text{ for some eigenvector $x$ of } (E+\Delta_E,J+\Delta_J) \}\nonumber\\
		&=\inf\{ \nnrm{(\Delta_E,\Delta_J,\Delta_R)}:(\Delta_E,\Delta_J,\Delta_R)\in \mathbb S, (R+\Delta_R)x=0 \text{ for some $x\in \C^n\setminus\{0\}$ satisfying } \nonumber\\ &\hspace{5cm}  (J+\Delta_J)x= i\omega (E+\Delta_E) x, \omega\in \R\}\nonumber\\
		&= \inf_{\omega\in \R}\inf_{x\in\C^{n}\setminus \{0\}} \Big(\inf \{ \nnrm{(\Delta_E,\Delta_J,\Delta_R)}:(\Delta_E,\Delta_J,\Delta_R)\in \mathbb S, (R+\Delta_R)x=0, (J+\Delta_J)x= i\omega (E+\Delta_E) x\} \Big)\nonumber\\
		&=\inf_{\omega \in \R } \rho_\omega^{\mathbb S},
    \end{align*}
    %
    % \begin{align}
    %     d_{im}^{\mathbb S}(E,J,R)&= \inf\{\nnrm{(\Delta_E,\Delta_J,\Delta_R)}~:~(\Delta_E,\Delta_J,\Delta_R)\in \mathbb S,~\exists~v_E,v_J\in \C^{n}, \text{ s.t. } v:=-i\omega v_E+v_J\neq 0,\nonumber\\
    %     & \hspace{5.5cm}\Delta_RM v= -RM v,~\Delta_EM v=v_E,~\Delta_JM v=v_J \label{eq:dist-im1} \}\\
    %     &= \inf_{\omega\in \R}\inf_{v_E,v_J}\{\nnrm{(\Delta_E,\Delta_J,\Delta_R)}~:~(\Delta_E,\Delta_J,\Delta_R)\in \mathbb S,~\exists~v_E,v_J\in \C^{n}, \text{ s.t. } v:=-i\omega v_E+v_J\neq 0,\nonumber\\
    %     & \hspace{5.5cm}\Delta_RM v= -RM v,~\Delta_EM v=v_E,~\Delta_JM v=v_J \label{eq:dist-im1} \}\\
    %     &= \inf_{\omega\in \R}\rho(\omega).
    % \end{align}
    %
    where, for a given scalar $\omega \in \R$, we have
    %
    \begin{align}\label{eq:rho}
        \rho_\omega^{\mathbb S}:=\inf_{x\in \C^{n}\setminus \{0\}}\Big(\inf\{\nnrm{(\Delta_E,\Delta_J,\Delta_R)}^2:(\Delta_E,\Delta_J,\Delta_R)\in \mathbb S,(R+\Delta_R)x=0, (J+\Delta_J)x= i\omega (E+\Delta_E) x\}\Big).
    \end{align}
    %
    By using the set $\Lambda$ of eigenvalues of the pair $(iE,J)$, we obtain
    %
    \begin{equation}\label{eq:dist-im1}
        d_{im}^{\mathbb S}(E,J,R)=\inf_{\omega\in \Lambda \cup \R\setminus\Lambda} \rho_{\omega}^{\mathbb S}.
    \end{equation}
    %
    We first consider the case when $\mathbb S=\mathcal S_d(E,J,R)$. Let $\omega\in \R\setminus\Lambda$, then, using Lemma~\ref{lem:tfae-map} in~\eqref{eq:rho}, we obtain
    %
    \begin{align}
        \rho_\omega^{\mathcal S_d}\Big|_{\R\setminus \Lambda}&= \inf_{x\in \C^{n}\setminus \{0\}}\Big(\inf\{\nnrm{(\Delta_E,\Delta_J,\Delta_R)}^2:(\Delta_E,\Delta_J,\Delta_R)\in S_d,\exists~v_E,v_J\in \C^{n}, \text{ s.t. } v:=-i\omega v_E+v_J\neq 0,\nonumber\\
        & \hspace{5.5cm}\Delta_RM v= -RM v,~\Delta_EM v=v_E,~\Delta_JM v=v_J\} \Big).\label{eq:rho1}
    \end{align}
    %
    To proceed, we solve the associated mapping problems under the imposed structure on the set $\mathcal S_d(E,J,R)$. By~{Theorem~\ref{map:def}}, there exist matrices $\Delta_E\leq 0,$ $\Delta_R\leq 0$ and by Theorem~\ref{map:herm} there exists skew-Hermitian $\Delta_J$ satisfying the constraints in \eqref{eq:rho1} if and only if
    %
    \begin{equation}
        v^*M^*RM v>0,~~ v_E^*M v>0,~ \text{  and }~ v_J^*M v \in i\R.
    \end{equation}
    %
    Moreover, among all admissible mappings, the minimal norms are given by
    %
    \begin{equation}
        \|\Delta_R\|=\frac{\|RMv\|^2}{v^*M^*RM v},~\|\Delta_E\|=\frac{\|v_E\|^2}{v_E^*M v},~\text{ and }~\|\Delta_J\|=\frac{\|v_J\|}{\|M v\|}.
    \end{equation}
    %
    Introducing the vectors $u=[v_J^T,~v_E^T]^T$, $L=[I_n,~-i\omega I_n]$, and substituting the minimal norm mappings in~\eqref{eq:rho1}, we obtain
    %
    \begin{align}
        \rho_\omega^{\mathcal S_d}\Big|_{\R\setminus \Lambda}\geq {\inf_{\substack{u\in\C^{2n},\\Lu\neq 0}}\Bigg\{\left(\frac{\|RMLu\|^2}{u^*L^*M^*RMLu}\right)^2+\frac{\|[I_n~0] u\|^2}{\|MLu\|^2} + \left(\frac{\|[0~I_n] u\|^2}{u^*L^*M^*[0~I_n] u}\right)^2: \begin{matrix} u^*[I_n~0]^TMLu \in i\R,\\ ~u^*L^*M^*[0~I_n]u>0 \end{matrix}\Bigg\}}.
        \label{rhobound}
    \end{align}
    %
    The inequality in \eqref{rhobound} arises because, in solving the individual minimal-norm mapping problems, we did not explicitly enforce the additional conditions $E+\Delta_E\geq 0$ and $R+\Delta_R\geq 0$.

    We now turn to the case $\omega \in \Lambda$. In this situation, the minimal norm $\nnrm{(\Delta_E,\Delta_J,\Delta_R)}$ is attained by choosing $\Delta_E=0$ and $\Delta_J=0$, and by selecting $\Delta_R$ such that $\Delta_R x_\omega = - R x_\omega$, where $x_\omega$ is an eigenvector of the matrix pair $(iE,J)$ corresponding to the eigenvalue $\omega$.
   Note that $x_\omega^*Rx_\omega \neq 0$. Indeed, if $x_\omega^*Rx_\omega = 0$ then $Rx_\omega=0$ as $R \geq 0$. By Lemma~\ref{lem:tfae-map} this implies that $i\omega$ is an eigenvalue of the dHDAE $(E,J-R)$ on the imaginary axis, which is a contradiction as the pair $(E,J-R)$ is assumed to be robustly asymptotically stable.

    Among all Hermitian negative semidefinite perturbations $\Delta_R \leq 0$ satisfying this constraint, the minimal norm is given by
    \[
    \|\Delta_R\| = \frac{\|R x_\omega\|^2}{x_\omega^* R x_\omega},
    \]
    and this minimum is attained for
    \[
    \Delta_R = -\frac{(R x_\omega)(R x_\omega)^*}{x_\omega^* R x_\omega}.
    \]
    Note that, in view of~\cite[Lemma 4.1]{MehMS16}, this choice of $\Delta_R$ also satisfies $R+\Delta_R \geq 0$. Substituting this expression into~\eqref{eq:rho}, we obtain
    \[
    \rho_\omega^{\mathcal S_d}\big|_{\Lambda} = \inf_{\omega \in \Lambda} \frac{{\|R x_\omega\|}^4}{{(x_\omega^* R x_\omega)}^2},
    \]
    where $x_\omega$ denotes eigenvector associated with the eigenvalue $\omega$. Combining the cases $\omega\in\mathbb{R}\setminus\Lambda$ and $\omega\in\Lambda$, we obtain the bound in~\eqref{Sd:dist-imagEJR}.\\

    We now consider the case $\mathbb S = \mathcal S_i(E,J,R)$. For this structured perturbation set, it follows from~\eqref{eq:rho1} that, for $\omega \in \mathbb{R}\setminus\Lambda$,
    %
    \begin{align}
        \rho_\omega^{\mathcal S_i}\Big|_{\R\setminus \Lambda}&= \inf_{x\in \C^{n}\setminus \{0\}}\Big(\inf\{\nnrm{(\Delta_E,\Delta_J,\Delta_R)}^2~:~(\Delta_E,\Delta_J,\Delta_R)\in \mathcal S_i,~\exists~v_E,v_J\in \C^{n}, \text{ s.t. } v:=-i\omega v_E+v_J\neq 0,\nonumber\\
        & \hspace{5.5cm}\Delta_RM v= -RM v,~\Delta_EM v=v_E,~\Delta_JM v=v_J\} \Big).\label{eq:rho3}
    \end{align}
    %
    We proceed analogously to the case $\mathbb S=\mathcal S_d(E,J,R)$ by solving the associated minimal-norm mapping problems
    %
    \begin{equation*}
        \Delta_R M v = - R M v, \quad
    \Delta_E M v = v_E, \quad
    \Delta_J M v = v_J,
    \end{equation*}
    %
    with $(\Delta_E,\Delta_J,\Delta_R)\in \mathcal S_i(E,J,R)$. By Theorem~\ref{map:herm}, there exist Hermitian matrices $\Delta_R,\Delta_E$ and a skew-Hermitian matrix $\Delta_J$ %\in\mathrm{SHerm}(n)$
    satisfying these constraints if and only if
    %
    \begin{equation*}
        v^* M^* R M v \in \mathbb{R}, \qquad v_E^* M v \in \mathbb{R}, \qquad v_J^* M v \in i\mathbb{R}.
    \end{equation*}
    %
    The first condition holds trivially since $R$ is Hermitian. The remaining two conditions can equivalently be written as
    %
    \begin{equation*}
        u^* \tilde H_1 u = 0, \quad u^* \tilde H_2 u = 0,
    \end{equation*}
    %
    where $u := [\, v_J^T \;\; v_E^T \,]^T$ and the matrices $\tilde H_1$ and
    $\tilde H_2$ are defined in~\eqref{mat:tildeH}.

    Among all feasible mappings, the minimal norms are given by
    %
    \begin{equation*}
        \|\Delta_R\| = \frac{\|R M v\|}{\|M v\|}, \quad
        \|\Delta_E\| = \frac{\|v_E\|}{\|M v\|}, \quad
        \|\Delta_J\| = \frac{\|v_J\|}{\|M v\|}.
    \end{equation*}
    %
    Hence,
    %
    \begin{equation*}
        \nnrm{(\Delta_E,\Delta_J,\Delta_R)}^2 = \frac{\|v_E\|^2 + \|v_J\|^2 + \|R M v\|^2}{\|M v\|^2} = \frac{u^* G_1 u}{u^* G_2 u},
    \end{equation*}
    %
    where the matrices $G_1$ and $G_2$ are defined in~\eqref{mat:G1G2}. Using these formulas in~\eqref{eq:rho3}, we obtain
    %
    \begin{equation}
        \rho_\omega^{\mathcal S_i}\Big|_{\R\setminus \Lambda} \geq \inf_{u\in \C^{2n}\setminus \{0\}} \left\{ \frac{u^* G_1 u}{u^* G_2 u}~:~ u^*G_2u\neq 0,~ u^*\tilde H_1u=0,~u^*\tilde H_2u=0 \right\}.\label{rhobound1}
    \end{equation}
    %
    Note that we have an inequality in \eqref{rhobound1} as we did not impose the conditions $R+\Delta_R\geq0, E+\Delta_E\geq 0$ from the set $\mathcal S_i(E,J,R)$ while solving the mapping problems. Finally, since $G_1$ is positive definite, we may introduce its principal positive definite square root and obtain the equivalent formulation
    %
    \begin{equation}\label{eq:temp1}
        \rho_\omega^{\mathcal S_i}\Big|_{\R\setminus \Lambda} \geq \left( \sup_{y\in \C^{2n}\setminus \{0\}} \left\{ \frac{y^* G y}{y^*y}~:~ y^*H_1y=0,~y^*H_2y=0 \right\}\right)^{-1},
    \end{equation}
    %
    where $y:=G_1^{1/2}u$, and the matrices $G:=G_1^{-1/2}G_2G_1^{-1/2},H_1:=G_1^{-1/2}\tilde H_1G_1^{-1/2},$ and $H_2:=G_1^{-1/2}\tilde H_2G_1^{-1/2}$ are as in~\eqref{mat:GH1H2}. Note that in~\eqref{eq:temp1} we omitted the condition $u^*G_2u\neq 0$ (or equivalently $y^*Gy\neq 0$), as the distance to a pencil with purely imaginary eigenvalues $d_{im}^{\mathcal S_i}(E,J,R)$ is finite, therefore including vectors $u$ such that $u^*G_2u=0$, will not effect the supremum.

Also note that $\alpha H_1+\beta H_2$ is indefinite (i.e., strictly not semidefinite) for every $(\alpha,\beta)\in \R^2 \setminus \{0\}$. Thus in view of~\cite[Theorem 3.2]{BorKMS14}, we have
%
\[
\rho_\omega^{\mathcal S_i}\Big|_{\R\setminus \Lambda}
\geq \left(\min_{t_1,t_2 \in \R} \lambda_{\max}(G+t_1H_1+t_2H_2) \right)^{-1}.
\]
%
By following the arguments similar to the case $\rho_\omega^{\mathcal S_d}\Big|_{ \Lambda}$ and using the minimal norm Hermitian mapping from Theorem~\ref{map:herm}, we obtain
%
  \[
    \rho_\omega^{\mathcal S_i}\big|_{\Lambda} = \inf_{\omega \in \Lambda} \frac{{\|R x_\omega\|}^2}{{\|x_\omega\|}^2},
    \]
    %
  where $x_\omega$ is the eigenvector associated with the eigenvalue $\omega$. Combining the cases $\omega\in\mathbb{R}\setminus\Lambda$ and $\omega\in\Lambda$, we obtain the bound in~\eqref{Si:dist-imagEJR}.
\end{proof}
%

\begin{remark}{\rm
    Note that Theorem~\ref{thm:dist-imagEJR} provides only a lower bound for the smallest perturbation that moves an eigenvalue to the imaginary axis. This is because, in solving the associated minimal-norm mapping problems in the proof, we did not enforce the additional constraints in the perturbation set $\mathbb S \in \{\mathcal S_d(E,J,R),\mathcal S_i(E,J,R)\}$, namely the conditions $E+\Delta_E \ge 0$ and $R+\Delta_R \ge 0$.

    If the infimum in~\eqref{Sd:dist-imagEJR} and~\eqref{Si:dist-imagEJR} is attained at some $(\hat \omega,\hat x)$, then the corresponding optimal perturbations $(\hat \Delta_E$, $\hat \Delta_J, \hat \Delta_R)$ can be constructed explicitly as described in the proof of Theorem~\ref{thm:dist-imagEJR}. Moreover, if these perturbations additionally satisfy $E+\hat \Delta_E \ge 0$ and $R+\hat \Delta_R \ge 0$, then the lower bound in~\eqref{Sd:dist-imagEJR} and~\eqref{Si:dist-imagEJR} are tight, and equality holds.
    }
\end{remark}
%
After characterizing the distance to the nearest dHDAE with purely imaginary eigenvalues, in the next subsection we treat the distance to the nearest singular dHDAE.
\subsection{Distance to the nearest singular dHDAE}\label{sec:distsing}

The unstructured distance to singularity $d_{sing}(E,J,R):=d_{sing}^{\mathcal S_d}{\mathbb(E,J,R)}$, when $\mathbb S={(\C^{n,n})}^3$ in~\eqref{def:dist-sing} and the structured distance $d_{sing}^{\mathcal S_d}(E,J,R)$ were studied in~\cite{PraS22}. For the sake of completeness, we state the following result from~\cite{PraS22} in our notation.

\begin{theorem}\label{thm:dist-singEJR-d}
        Consider an asymptotically stable dHDAE system of the form~\eqref{dhdae1}. Then
    %
    \begin{enumerate}
        \item the unstructured distance to singularity $d_{sing}(E,J,R)$ is given by
    %
    \begin{equation}\label{Si:dist-sing-EJR1}
        d_{sing}(E,J,R)=\sqrt{\lambda_{\min}(E^2-J^2+R^2)}.
    \end{equation}
    %
        \item the distance to singularity $d_{sing}^{\mathcal S_d}(E,J,R)$ with respect to perturbations from the set $\mathcal S_d(E,J,R)$ is given by
\begin{equation}\label{Sd:dist-sing-EJR}
(d_{sing}^{\mathcal S_d}(E,J,R))^2=\min\left\{d_{sing}^{\mathcal S_d}(J,R)^2,d_{sing}^{\mathcal S_d}(E,J)^2,
\min_{\alpha \in \mathcal M_3} \frac{\alpha^*J^*J\alpha}{\alpha^*\alpha}+
 \left(\frac{\alpha^*R^2\alpha}{\alpha^*R\alpha}\right)^2+ \left(\frac{\alpha^*E^2\alpha}{\alpha^*E\alpha}\right)^2
\right\},
   \end{equation}
where  $\mathcal M_3=\text{\rm ker}(R)^{c}\cap \text{\rm ker}(E)^{c}$,
and $d_{sing}^{\mathcal S_d}(J,R)$ and
$d_{sing}^{\mathcal S_d}(E,J)$ are given by~\cite[Table A.1]{PraS22}.
    \end{enumerate}
\end{theorem}


\begin{theorem}\label{thm:dist_singEJR_i}
    Consider an asymptotically stable dHDAE system of the form~\eqref{dhdae1}. Then
the distance to singularity $d_{sing}^{\mathcal S_i}(E,J,R)$ with respect to perturbations from the set $\mathcal S_i(E,J,R)$ is bounded by
    %
    \begin{equation}\label{Si:dist-sing-EJR}
    d_{sing}^{\mathcal S_d}(E,J,R) \geq    d_{sing}^{\mathcal S_i}(E,J,R)\geq \sqrt{\lambda_{\min}(E^2-J^2+R^2)}.
    \end{equation}
    Furthermore, if $R>0$ and $E>0$, then
    \begin{equation}\label{Si:dist-sing-EJR2}
   d_{sing}^{\mathcal S_i}(E,J,R)= \sqrt{\lambda_{\min}(E^2-J^2+R^2)}= d_{sing}(E,J,R).
    \end{equation}
    %
\end{theorem}
\begin{proof} The inequalities in~\eqref{Si:dist-sing-EJR} are immediate by Definition~\ref{def:distances}, since
%
\[
S_d(E,J,R) \subseteq S_i(E,J,R) \subseteq (\C^{n,n})^3.
\]
%
In view of~\eqref{def:dist-sing} and Corollary~\ref{cor}, we have
	\begin{align}
		&d_{sing}^{\mathcal S_i}(E,J,R)\nonumber\\
		&=\inf \left\{ \nnrm{(\Delta_E,\Delta_J,\Delta_R)}:(\Delta_E,\Delta_J,\Delta_R)\in \mathcal S_i,~ \ker(E+\Delta_E) \cap \ker(J + \Delta_J) \cap \ker(R+\Delta_R) \neq \{0\} \right\}\nonumber\\
		&=\inf \{ \nnrm{(\Delta_E,\Delta_J,\Delta_R)}:(\Delta_E,\Delta_J,\Delta_R)\in \mathcal S_i,~x\in \C^{n}\setminus\{0\},~ (E+\Delta_E)x=0, (J + \Delta_J)x=0,\nonumber \\
        &\hspace{10cm}(R+\Delta_R)x=0 \}.\label{Sd:dist-sing-EJR1}
	\end{align}
	%
For a given $x\in \C^{n}\setminus \{0\}$, to obtain the minimal norm perturbation $(\Delta_E,\Delta_J,\Delta_R) \in \mathcal S_i(E,J,R)$ we have to solve the mapping problem
	%
	\begin{equation}\label{mapping:constraints}
		(E+\Delta_E)x=0,~ \Delta_Rx=-Rx,~\text{and }~ \Delta_Jx=-Jx.
	\end{equation}
	%
    By Theorem~\ref{map:herm}, there exist Hermitian matrices $\Delta_R,\Delta_E$ and a skew-Hermitian matrix $\Delta_J$ satisfying these constraints if and only if
    %
    \begin{equation*}
        x^*R x \in \mathbb{R}, \qquad x^* E x \in \mathbb{R}, \qquad x^* J x \in i\mathbb{R},
    \end{equation*}
    %
    which holds trivially due to the structure of the matrices $E,J$ and $R$. Among all feasible mappings, the minimal norms are given by
    %
    \begin{equation*}
        \|\Delta_R\| = \frac{\|R x\|}{\|x\|}, \quad
        \|\Delta_E\| = \frac{\|Ex\|}{\|x\|}, \quad
        \|\Delta_J\| = \frac{\|Jx\|}{\|x\|}.
    \end{equation*}
    %
    Also when $E>0$ and $R>0$, from~\cite[Lemma 4.4]{MehMS16}, the minimal norm mappings $\Delta_E$ and $\Delta_R$ satify that $R+\Delta_R \geq 0$ and $E+\Delta_E \geq 0$.
    Using the above formulas in~\eqref{Sd:dist-sing-EJR1}, we obtain
    %
    \begin{equation*}
        (d_{sing}^{\mathcal S_i}(E,J,R))^2=\inf_{x\in \C^{n}\setminus\{0\}} \frac{x^*(E^*E+J^*J+R^*R)x}{x^*x} = \lambda_{\min}(E^2-J^2+R^2).
    \end{equation*}
    This completes the proof.
\end{proof}


\subsection{Distance to the nearest dHDAE with higher index}\label{sec:disthi}
Having established an explicit characterization of the distance to singularity for dHDAE systems,
%under full structured perturbations,
we now turn our attention to the nearest system of higher index
%. In the following, we therefore investigate the distance to higher index
under the same classes of structure-preserving perturbations, $ {\mathcal S_d}(E,J,R)$ and ${\mathcal S_i}(E,J,R).$

\begin{theorem}\label{thm:dist-hiEJR}
    Consider a robustly asymptotically stable dHDAE system of the form~\eqref{dhdae1}. Then
    %
    \begin{enumerate}
        \item the distance to higher index $d_{hi}^{\mathcal S_d}(E,J,R)$ with respect to perturbations from the set $\mathcal S_{d}(E,J,R)$ is bounded via
        %
        \begin{equation}
            (d_{hi}^{\mathcal S_d}(E,J,R))^2 \leq \min_{1 \leq k \leq n} \inf_{x\neq 0} \left\{ \lambda_{n-k+1}(E)^2 + \frac{\|(N_k^*)^{\dagger} N_k^* J N_k x\|^2}{\|N_k x\|^2} + \frac{\|(N_k^*)^{\dagger} N_k^* R N_k x\|^4} {\big(x^* N_k^* R N_k x\big)^2} \right\},
        \end{equation}
        %
        where $N_k$ denotes a  matrix with  columns that span a basis of $\ker(E+\hat \Delta_E)$, $\hat \Delta_E$ is the optimal perturbation such that $\rank (E+\hat \Delta_E)=k$,
         and $\lambda_{n-k+1}(E)$ is the $(n-k+1)$th smallest eigenvalue of $E$.
        \item if $R>0$, the distance to higher index $d_{hi}^{\mathcal S_i}(E,J,R)$ with respect to perturbations from the set $\mathcal S_{i}(E,J,R)$ is bounded via
        %
        \begin{equation}
            (d_{hi}^{\mathcal S_i}(E,J,R))^2 \leq \min_{1\leq k \leq n}\inf_{x\neq 0} \left\{ \lambda_{n-k+1}(E)^2 + \frac{\|(N_k^*)^{\dagger} N_k^* J N_k x\|^2}{\|N_k x\|^2} + \frac{\|(N_k^*)^{\dagger} N_k^* R N_k x\|^2} {\|N_k x\|^2} \right\},
        \end{equation}
        %
        where $N_k$ denotes a  matrix with  columns that span a basis of $\ker(E+\hat \Delta_E)$, $\hat \Delta_E$ is the optimal perturbation such that $\rank (E+\hat \Delta_E)=k$, and $\lambda_{n-k+1}(E)$ is the $(n-k+1)$th smallest eigenvalue of $E$.
    \end{enumerate}
\end{theorem}
\begin{proof}
    Using Definition~\eqref{def:distances}, the distance to higher index under structured perturbations from the set $\mathbb S\in \{\mathcal S_d,\mathcal S_i\}$ is given by
    \[
    d_{hi}^{\mathbb S}(E,J,R) = \inf\Big\{\nnrm{(\Delta_E,\Delta_J,\Delta_R)} ~:~ (\Delta_E,\Delta_J,\Delta_R)\in \mathbb S,~ \text{ index of } (E+\Delta_E,J+\Delta_J-R-\Delta_R) > 1 \Big\}.
    \]
    In view of Theorem~\ref{thm:stabledh}, this can be equivalently written as
    %
    \begin{align*}
        &d_{hi}^{\mathbb S}(E,J,R) \\ &= \inf\Big\{ \nnrm{(\Delta_E,\Delta_J,\Delta_R)} ~:~ (\Delta_E,\Delta_J,\Delta_R)\in \mathbb S,~ \text{rank}\big(N(E+\Delta_E)^* (J+\Delta_J-R-\Delta_R) N(E+\Delta_E) \big) < n \Big\},
    \end{align*}
    %
    where $N(E+\Delta_E)$ denotes a basis matrix for the null space of $E+\Delta_E$.

    Observe that the constraint in the above formulation involves second-order perturbation terms through products of $\Delta_E$ with $\Delta_J$ and $\Delta_R$. To address this difficulty, we adopt a two-stage approach. First, we determine an optimal perturbation $\Delta_E$ such that $\dim \ker(E+\Delta_E) = k$. Then, for this fixed $\Delta_E$, we determine the minimal structured perturbations $(\Delta_J,\Delta_R)$ that satisfy the rank condition. Let $E = U \Sigma U^*$ be the spectral decomposition of $E$, with $\Sigma = \diag \{\lambda_1,\lambda_2,\ldots,\lambda_n\}$, where $\lambda_1 \ge \cdots \ge \lambda_n$. The optimal perturbation $\Delta_E$ achieving ${\rank}(E+\Delta_E)=n-k$ is given by
    \[
    \hat \Delta_E = - U \diag\{0,0,\ldots,0,\lambda_{n-k+1},\lambda_{n-k+2}, \ldots,\lambda_{n-1},\lambda_{n}\}U^*.% - \lambda_n x_n x_n^* - \lambda_{n-1} x_{n-1} x_{n-1}^* - \cdots - \lambda_{n-k+1} x_{n-k+1} x_{n-k+1}^*,
    \]
    %where $x_{n-k+1}, \ldots,x_n$ are the corresponding right eigenvectors of $E$.
    Let $N_k$ be a  matrix whose columns form a basis for $\ker(E+\hat \Delta_E)$.
    %With this choice of $\hat \Delta_E$,
    Then we obtain
    %
    \begin{align}\label{Sd:muintro}
        (d_{hi}^{\mathbb S}(E,J,R))^2 &\leq  \inf\Big\{ \lambda_{n-k+1}(E)^2 + \|\Delta_J\|^2 + \|\Delta_R\|^2 : (\hat \Delta_E,\Delta_j,\Delta_R)\in \mathbb S,\nonumber \\ & \hspace{3cm} \exists\, x\neq 0 \text{ such that } N_k^* (J+\Delta_J-R-\Delta_R) N_k x = 0 \Big\}
        \nonumber\\
        &=\lambda_{n-k+1}(E)^2 + \mu_k^{\mathbb S},
    \end{align}
    where
    %
    \begin{eqnarray}\label{hu_hi}
        \mu_k^{\mathbb S}:=\inf\Big\{ \|\Delta_J\|^2 + \|\Delta_R\|^2 : (\hat \Delta_E,\Delta_j,\Delta_R)\in \mathbb S,\exists\, x\neq 0 \text{ such that } N_k^* (J+\Delta_J-R-\Delta_R) N_k x = 0 \Big\}.
    \end{eqnarray}
    %
    Since $J+\Delta_J$ is skew-Hermitian and $R+\Delta_R$ is Hermitian positive semidefinite, the condition
    \[
    (J+\Delta_J - R - \Delta_R)y = 0
    \]
    holds if and only if
    \[
    (J+\Delta_J)y = 0 \quad \text{and} \quad (R+\Delta_R)y = 0.
    \]
    Therefore,
    %
    \begin{eqnarray}\label{eq:temp2}
        \mu_k^{\mathbb S} &= & \inf\Big\{  \|\Delta_J\|^2 + \|\Delta_R\|^2 : (\hat \Delta_E,\Delta_j,\Delta_R)\in \mathbb S, \exists\, x\neq 0 \text{ such that } \nonumber
        \\ && \hspace{4.5cm} N_k^* (J+\Delta_J) N_k x = 0,~ N_k^* (R+\Delta_R) N_k x = 0 \Big\}.
    \end{eqnarray}
    %
    In~\cite[Lemma 2.8]{MehMS16} it is shown that for a matrix $B$, the equation $B\Delta x = y$ holds if and only if $\Delta x = B^{\dagger}y$ and $BB^{\dagger}y = y$, where $B^{\dagger}$ denotes the Moore-Penrose pseudoinverse. Applying this to the above constraints in~\eqref{eq:temp2}, the relations
    \[
    N_k^* \Delta_J N_k x = - N_k^* J N_k x, \qquad N_k^* \Delta_R N_k x = - N_k^* R N_k x
    \]
    can be equivalently written as
    \[
    \Delta_J N_k x = - (N_k^*)^{\dagger} N_k^* J N_k x, \qquad \Delta_R N_k x = - (N_k^*)^{\dagger} N_k^* R N_k x,
    \]
    since the consistency conditions are trivially satisfied. Hence, we obtain
    %
    \begin{eqnarray}\label{Sd:high1}
        \mu_k^{\mathbb S} &=  \inf\Big\{  \|\Delta_J\|^2 + \|\Delta_R\|^2 : (\hat \Delta_E,\Delta_J,\Delta_R)\in \mathbb S,~\exists\, x\neq 0, \; \Delta_J N_k x = - (N_k^*)^{\dagger} N_k^* J N_k x, \;
        \nonumber \\ &\hspace {5cm}\Delta_R N_k x = - (N_k^*)^{\dagger} N_k^* R N_k x \Big\}.
    \end{eqnarray}
    %
    We now solve the associated minimal-norm mapping problems.
    First, let $\mathbb S=\mathcal S_d(E,J,R)$. By Theorem~\ref{map:herm} and~\cite[Theorem~2.3]{MehMS16}, there exist a skew-Hermitian matrix $\Delta_J$ and a Hermitian negative semidefinite matrix $\Delta_R\le 0$ satisfying the constraints in~\eqref{Sd:high1} if and only if
    \[
    x^* N_k^* J N_k x \in i\mathbb{R}
    \quad \text{and} \quad
    x^* N_k^* R N_k x > 0.
    \]
    These conditions are automatically satisfied due to the structural properties of the matrices $J$ and $R$.
    %
    Among all such admissible mappings, the minimal norms are given by
    \[
    \|\Delta_J\| = \frac{\|(N_k^*)^{\dagger} N_k^* J N_k x\|}{\|N_k x\|}, \qquad
    \|\Delta_R\| = \frac{\|(N_k^*)^{\dagger} N_k^* R N_k x\|^2}{x^* N_k^* R N_k x},
    \]
    and are attained by
    \[
    \Delta_J = i\,\hat H_{(N_k x,i(N_k^*)^{\dagger} N_k^* J N_k x)}, \qquad
    \Delta_R = -\frac{\big((N_k^*)^{\dagger} N_k^* R N_k x\big) \big((N_k^*)^{\dagger} N_k^* R N_k x\big)^*}{x^* N_k^* R N_k x},
    \]
    where $\hat H$ is defined in~\eqref{def:hatH}.
    %
    If $R\ge 0$ is singular and $x^* N_k^* R N_k x = 0$, then necessarily $(N_k^*)^{\dagger} N_k^* R N_k x = 0$. In this case, adopting the convention $0/0 = 0$, the choice $\Delta_R = 0$ satisfies both the mapping constraint in~\eqref{Sd:high1} and the minimal-norm condition. Moreover, the resulting perturbation satisfies $R + \Delta_R \ge 0$; see~\cite[Lemma~4.1]{MehMS16}.
    %
    Combining these observations in~\eqref{Sd:high1}, we obtain
    \begin{equation}\label{Sd:high11}
    \mu_k^{\mathcal S_d} =  \inf_{x \neq 0} \left\{
    \frac{\|(N_k^*)^{\dagger} N_k^* J N_k x\|^2}{\|N_k x\|^2} +
    \frac{\|(N_k^*)^{\dagger} N_k^* R N_k x\|^4}{\big(x^* N_k^* R N_k x\big)^2} \right\},
    \end{equation}
       and thus from~\eqref{Sd:muintro} and~\eqref{Sd:high11}, we have
    \[
    (d_{hi}^{\mathcal S_d}(E,J,R))^2 \le \min_{1 \le k \le n} \inf_{x \neq 0} \left\{ \lambda_{n-k+1}(E)^2 +
    \frac{\|(N_k^*)^{\dagger} N_k^* J N_k x\|^2}{\|N_k x\|^2} +
    \frac{\|(N_k^*)^{\dagger} N_k^* R N_k x\|^4}{\big(x^* N_k^* R N_k x\big)^2} \right\}.
    \]
    Next, consider $\mathbb S=\mathcal S_i(E,J,R)$. By Theorem~\ref{map:herm}, there exist a skew-Hermitian matrix $\Delta_J$ and a Hermitian matrix $\Delta_R$ satisfying the constraints in~\eqref{Sd:high1} if and only if
    \[
    x^* N_k^* J N_k x \in i\mathbb{R}
    \quad \text{and} \quad
    x^* N_k^* R N_k x \in \mathbb{R},
    \]
    which again holds trivially due to the structure of $J$ and $R$.
    %
    The corresponding minimal norms are
    \[
    \|\Delta_J\| = \frac{\|(N_k^*)^{\dagger} N_k^* J N_k x\|}{\|N_k x\|},
    \qquad
    \|\Delta_R\| = \frac{\|(N_k^*)^{\dagger} N_k^* R N_k x\|}{\|N_k x\|},
    \]
    and are attained by
    \[
    \Delta_J = -i\,\hat H_{(N_k x,\, -i (N_k^*)^{\dagger} N_k^* J N_k x)},
    \qquad
    \Delta_R = \hat H_{(N_k x,\,-(N_k^*)^{\dagger} N_k^* R N_k x)},
    \]
    where $\hat H$ is defined in~\eqref{def:hatH}. Note that $(R+\Delta_R)Nx=0$, hence on applying~\cite[Lemma 4.4]{MehMS16} we have $R+\Delta_R\geq 0$.
    %
    Applying these minimal-norm mappings to~\eqref{Sd:high1} yields
    \begin{equation}\label{Sd:high22}
    \mu_k^{\mathcal S_i} =  \inf_{x \neq 0} \left\{
    \frac{\|(N_k^*)^{\dagger} N_k^* J N_k x\|^2}{\|N_k x\|^2} +
    \frac{\|(N_k^*)^{\dagger} N_k^* R N_k x\|^2}{\|N_k x\|^2} \right\}.
    \end{equation}
    Thus in view of~\eqref{Sd:muintro} and~\eqref{Sd:high22}, we have
    \[
    (d_{hi}^{\mathcal S_i}(E,J,R))^2 \le \min_{1 \le k \le n} \inf_{x \neq 0} \left\{
    \lambda_{n-k+1}(E)^2 +
    \frac{\|(N_k^*)^{\dagger} N_k^* J N_k x\|^2}{\|N_k x\|^2} +
    \frac{\|(N_k^*)^{\dagger} N_k^* R N_k x\|^2}{\|N_k x\|^2} \right\}.
    \]
    This completes the proof.
    %
    % For $\mathbb S=\mathcal S_d(E,J,R)$, the minimal-norm structured mappings can be obtained from Theorem~\ref{map:herm} for optimal $\Delta_J^*=-\Delta_J$ and from~{Theorem~\ref{map:def}} for optimal $\Delta_R\leq 0$.d  The existense condition for these mappings are trivially satisfied. Using the minimal norm mappings we obtain
    % \[
    % (d_{hi}^{\mathcal S_d}(E,J,R))^2 \leq \min_{1\leq k \leq n}\inf_{x\neq 0} \left\{ \lambda_{n-k+1}(E)^2 + \frac{\|(N_k^*)^{\dagger} N_k^* J N_k x\|^2}{\|N_k x\|^2} + \frac{\|(N_k^*)^{\dagger} N_k^* R N_k x\|^4} {\big(x^* N_k^* R N_k x\big)^2} \right\}.
    % \]
    % In contrast, for $\mathbb S=\mathcal S_i(E,J,R)$, the corresponding minimal-norm mappings are characterized by Theorem~\ref{map:herm}, leading to
    % \[
    % (d_{hi}^{\mathcal S_i}(E,J,R))^2 \leq \min_{1\leq k \leq n}\inf_{x\neq 0} \left\{ \lambda_{n-k+1}(E)_n^2 + \frac{\|(N_k^*)^{\dagger} N_k^* J N_k x\|^2}{\|N_k x\|^2} + \frac{\|(N_k^*)^{\dagger} N_k^* R N_k x\|^2} {\|N_k x\|^2} \right\}.
    % \]
\end{proof}
%

\begin{remark}{\rm
    Observe that the definition of the distance to higher index $d_{hi}^{\mathbb S}(E,J,R)$ involves second-order perturbation terms, since the rank condition depends on products of perturbations. To overcome this difficulty, in the proof of the theorem we have employed a two-stage  procedure: first perturbing $E$ to increase the dimension of its null space, and then determining the minimal structured perturbations in $J$ and $R$. As a consequence, the result obtained in the above theorem provides only an upper bound for the distance to the nearest dHDAE system of higher index.

    At present, we cannot make a definitive statement about the tightness of this bound; nevertheless, it is expected to provide a good estimate in practice. In contrast, when perturbations are restricted to $J$ and $R$, i.e., in the computation of $d_{hi}^{\mathbb S}(J,R)$, then no second-order perturbation terms arise. This allows for an exact characterization of the distance to the nearest dHDAE of higher index, see the following section.
    }
\end{remark}

In this section we have derived characterizations and bounds for the three distances under structured perturbations to all three coefficients. Since in some applications the matrix $E$ is fixed and not subject to perturbations, in the next section we present perturbation results where the perturbation is restricted  to $\Delta_E=0$.


\section{Distances under  partial perturbations}\label{sec:partial}
%\subsection{Case 1: perturb only J and R}
In this section, we present analytic results for the various distances defined in Definition~\ref{def:distances} for dHDAE systems of the form~\eqref{dhdae1} by fixing $E$ and allowing only structured perturbations to $J$ and $R$, i.e. by considering the perturbation sets  as $\mathcal S_d(J,R):=\mathcal S_d(0,J,R)$
and $\mathcal S_i(J,R):=\mathcal S_i(0,J,R)$ in~\eqref{set:Sd} and~\eqref{set:Si}, respectively. The corresponding distances are denoted by $d_{sing}^{\mathbb S}(J,R)$, $d_{hi}^{\mathbb S}(J,R)$, $d_{im}^{\mathbb S}(J,R)$, and $d_{inst}^{\mathbb S}(J,R)$,
%and defined by~\eqref{def:dist-sing}--\eqref{def:dist-inst},
respectively, where $\mathbb S\in \{\mathcal S_d,\mathcal S_i\}$.

%{\color{red}
%The following two results Theorem 25 and 26 can %be mentioned as a remark to Theorem 27.
%}

If the matrix $E$ is nonsingular, then the distance to the boundary of the robustly asymptotically stable systems is equal to the distance to a dHDAE with purely imaginary eigenvalues, as the perturbed system will always be regular and of index zero. This distance was obtained in~\cite{BenMPS25} by transforming the dHDAE~\eqref{dhdae1} to an ordinary dH system of the form $\dot{x}=(J-R)Q$, where $Q=E^{-1}$.

% \begin{theorem}\label{thm:dist-im}
% 	Consider an asymptotically stable dHDAE system of the form~\eqref{dhdae1} with $E>0$. Then
% 	%
% 	\begin{enumerate}
% 		\item the distance to dH pencil with purely imaginary eigenvalues $d_{im}^{\mathcal S_d}(J,R)$ is given by
% 		\begin{equation}\label{eq:dist-im}
% 			(d_{im}^{\mathcal S_d}(J,R))^2= \inf_{\omega \in \R} \inf_{x\in \C^{n}\setminus \{0\}} \left\{\left(\frac{x^*R^2x}{x^*Rx} \right)^2+\frac{x^*(i\omega E -J)^*(i\omega E -J)x}{x^*x} \right\},
% 		\end{equation}
% 		\item if $R>0$, the distance to dH pencil with purely imaginary eigenvalues $d_{im}^{\mathcal S_i}(J,R)$ is given by
% 		\begin{equation}
% 			d_{im}^{\mathcal S_i}(J,R)=\inf_{\omega\in \R} \sigma_{\min}\left( \mat{c}R\\i\omega E - J \rix \right)
% 		\end{equation}
% 		and, if $R\geq 0$ is singular, then
% 		\begin{equation}
% 			d_{im}^{\mathcal S_i}(J,R)\geq \inf_{\omega\in \R} \sigma_{\min}\left( \mat{c}R\\i\omega E - J \rix \right).
% 		\end{equation}
% 	\end{enumerate}
% \end{theorem}
% \begin{proof}
% 	By Definition~\ref{def:distances}, we have
% 	%
% 	\begin{equation}
% 		d_{im}^{\mathcal S_d}(J,R)=\inf \left\{ \nnrm{(\Delta_J,\Delta_R)}~:~(\Delta_J,\Delta_R)\in\mathcal S_d(J,R),~ \Lambda(E,J+\Delta_J - (R+\Delta_R)) \cap i\R \neq \emptyset \right\}.
% 	\end{equation}
% 	%
% 	Since for $(\Delta_J,\Delta_R)\in S_d(J,R)$, the perturbed DH pair $(E,J+\Delta_J - (R+\Delta_R))$ still has the DH structure. By using~\cite[Lemma 4.1]{MehMS16}, we obtain
% 	%
% 	\begin{align}%\label{eq:def1}
% 		&d_{im}^{\mathcal S_d}(J,R) \nonumber\\
% 		&=\inf \{ \nnrm{(\Delta_J,\Delta_R)}:(\Delta_J,\Delta_R)\in\mathcal S_d(J,R),~ (R+\Delta_R)x=0 \text{ for some eigenvector $x$ of } (E,J+\Delta_J) \}\nonumber\\
% 		&=\inf\{ \nnrm{(\Delta_J,\Delta_R)}:(\Delta_J,\Delta_R)\in\mathcal S_d(J,R),~ (R+\Delta_R)x=0 \text{ for some $x\in \C^n\setminus\{0\}$ satisfying } \nonumber\\ &\hspace{5cm}  (J+\Delta_J)x= i\omega E x, \omega\in \R\}\nonumber\\
% 		&= \inf_{\omega\in \R}\inf_{x\in\C^{n}\setminus \{0\}} \Big(\inf \{ \nnrm{(\Delta_J,\Delta_R)}:(\Delta_J,\Delta_R)\in\mathcal S_d(J,R),~ (R+\Delta_R)x=0, (J+\Delta_J)x= i\omega E x\} \Big)\nonumber\\
% 		&=\inf_{\omega \in \R } \vartheta^{\mathcal S_d}_{\omega} \label{eq:def11}.
% 	\end{align}
% 	%
% 	%    where
% 	%the equation~\eqref{eq:def11} separated the variable $x$, $%\omega$ and $(\Delta_J,\Delta_R)$ into sublevel optimizations.
% 	For a given scalar $\omega$, we have
% 	%
% 	\begin{align}\label{eq:def12}
% 		\vartheta^{\mathcal S_d}_{\omega}&:= \inf_{x\in\C^{n}\setminus \{0\}} \Big(\inf \{ \nnrm{(\Delta_J,\Delta_R)}:(\Delta_J,\Delta_R)\in\mathcal S_d(J,R),~ (R+\Delta_R)x=0, \nonumber \\ & \hspace{6cm}(J+\Delta_J)x= i\omega E x\} \Big).
% 	\end{align}
% 	%
% 	For the inner optimization in~\eqref{eq:def12}, we will show now that the minimal value can be expressed as the sum of two generalized Rayleigh quotients, which depend on the variables $x$ and $\omega$. For this, we first solve two mapping problems.	
% 	For given $x\in \C^n\setminus\{0\}$ and $\omega\in \mathbb R$, determine the minimum norm solution $(\Delta_J,\Delta_R)$ such that $\Delta_Rx=-Rx$ and $\Delta_Jx=(i\omega E - J)x$. Using \cite[Theorem~\ref{map:herm}]{MacMT08}, there exists a skew-Hermitian $\Delta_J$ satisfying $\Delta_Jx=(i\omega E - J)x$ if and only if $\Re(x^*(i\omega E - J)x)=0$, which trivially holds because of the structure of $J$ and $E$. The minimal norm among all such mappings $\Delta_J$ is given by
% 	%
% 	\begin{equation}\label{norm:jdef}
% 		\|\Delta_J\|=\frac{\|(i\omega E - J)x\|}{\|x\|},
% 	\end{equation}
% 	%
% 	and is attained by $\Delta_J=i \hat H_{(Qx, (i\omega I_n-JQ)x)}$, where $\hat H$ is defined in~\eqref{def:hatH}. 	
% 	For the other constraint $\Delta_Rx=-Rx$. If $Rx\neq 0$, then from {Theorem~\ref{map:def}}, there exists $\Delta_R\leq 0$ satisfying $\Delta_Rx=-Rx$ if and only if $-x^*Rx<0$, which holds as $R\geq 0$, and the minimum norm solution %$\Delta_R$
% 	is given by
% 	%
% 	\begin{equation}\label{norm:rdef}
% 		\|\Delta_R\|=\frac{\|Rx\|^2}{x^*Rx},
% 	\end{equation}
% 	%
% 	which is attained by $\Delta_R=\frac{-1}{x^*Rx}(Rx)(Rx)^*$. If $Rx=0$ (this case may arise if $R$ is singular), then again from {Theorem~\ref{map:def}}, $\Delta_R=0$ is the solution that satisfies~\eqref{norm:rdef} by defining the fraction $\frac{0}{0}$ to have value $0$. Note that while solving the mapping constraint $\Delta_Rx=-Rx$, we have not considered the additional condition $R+\Delta_R\geq 0$. But using the fact that $S_d(J,R)\subseteq \{(\Delta_J,\Delta_R): \Delta_J^*=-\Delta_J, \Delta_R^*=\Delta_R\leq 0\}$, and using~\eqref{norm:jdef} and~\eqref{norm:rdef} in~\eqref{eq:def12}, we obtain
% 	%
% 	%	Similarly for $\Delta_R$, on using Lemma~\ref{map:def}, there exists $\Delta_R<0$ satisfying $\Delta_RQx=-RQx$ if and only if $-x^*QRQx<0$, which hold always because $R$ is Hermitian positive definite matrix. The minimal norm mapping $\Delta_R$ is given by
% 	%	%
% 	%	\begin{equation}\label{norm:rdef}
% 		%		\|\Delta_R\|=\frac{\|RQx\|^2}{x^*QRQx}
% 		%	\end{equation}
% 	%	%
% 	%	which is attained by $\Delta_R=\frac{1}{x^*QRQx}(RQx)(RQx)^*$.
% 	%	%
% 	%	Using~\eqref{norm:jdef} and~\eqref{norm:rdef} and the fact that $S_d(J,R)\subseteq \{(\Delta_J,\Delta_R): \Delta_J^*=-\Delta_J, \Delta_R^*=\Delta_R<0\}$, we obtain
% 	%	
% 	\begin{equation}\label{eq:def2}
% 		(\vartheta^{S_d}_{\omega})^2\geq \inf_{x\in\C^{n}\setminus\{0\}} \left\{ \frac{\|Rx\|^4}{(x^*Rx)^2}+\frac{\|(i\omega E - J)x\|^2}{\|x\|^2} \right\}.
% 	\end{equation}
% 	%
% 	If we can show that equality holds in~\eqref{eq:def2}, then the proof is finished. Note that the vector $x$ can be scaled to have unit norm without affecting the objective function in~\eqref{eq:def2}. Using the compactness of the set of unit norm vectors, the infimum in~\eqref{eq:def2} is actually a minimum and is attained for some $x=\hat x$. For a given $\omega$, set
% 	%
% 	\begin{equation}
% 		\hat{\Delta}_R=\begin{cases}
% 			-\frac{(R\hat x)(R\hat x)^*}{\hat x^*R\hat x} & \text{ if } R\hat{x} \neq 0,\\
% 			\qquad 0 & \text{ if } R\hat{x} = 0,
% 		\end{cases} \quad \quad \hat \Delta_J = \hat H_{(\hat x,(i\omega E - J)\hat x)},
% 	\end{equation}
% 	%
% 	where $\hat H$ is defined in~\eqref{def:hatH}. By {\cite[Lemma 4.1]{MehMS16}}, it then follows that for the optimal perturbation $\hat \Delta_R$ we have that $R+\hat \Delta_R\geq0$ so that equality holds in~\eqref{eq:def2}, as then $(\hat \Delta_J,\hat \Delta_R)\in S_d(J,R)$.
% 	%This follows directly by {\cite[Lemma 4.1]{MehMS16}} %on %using Lemma~\ref{lemdh4}. This completes the proof.
	
% 	{\color{red} Proof can be shortened or removed by citing the previous paper. Proof of the second part $d_{im}^{S_i}$ also follows on the similar lines.
%     This case can be transforemed to the standard DH system in the form $(J-R)Q$, where $Q=E^{-1}$.}
% \end{proof}
%
%\noindent It is straightforward to observe that
%
%\subsubsection{Bound for $r^{S_d}(J,R)$}
%Here, we discuss on some easily computable bounds for $r^{S_d}(J,R)$.
%It is easy to verify the following relation between %$r(J,R),r^{S}(J,R),r^{S_i}(J,R)$ and $r^{S_d}(J,R)$
%
% \begin{equation*}
% 	d_{im}^{\mathcal S_i}(J,R) \leq d_{im}^{\mathcal S_d}(J,R),
% \end{equation*}
% %
% since $\mathcal S_d(J,R)\subseteq \mathcal S_i(J,R)$. Thus, $d_{im}^{\mathcal S_i}(J,R)$ gives a lower bounds for $d_{im}^{\mathcal S_d}(J,R)$.
% Another lower bound for $d_{im}^{\mathcal S_d}(J,R)$
% was also obtained in~\cite[Corollary 3.2]{BenMPS25}.
% as follows,
%%
%\begin{align*}
%	r^{S_d}(J,R) & = \inf_{\omega \in \R} \inf_{x\in \C^{n}\setminus \{0\}} \left\{\left(\frac{x^*QR^2Qx}{x^*QRQx} \right)^2+\frac{x^*(i\omega I_n -JQ)^*(i\omega I_n -JQ)x}{x^*Q^2x} \right\} \\
%	& \geq \inf_{\omega \in \R} \inf_{x\in \C^{n}\setminus \{0\}} \left(\frac{x^*QR^2Qx}{x^*QRQx} \right)^2 + \inf_{\omega \in \R} \inf_{x\in \C^{n}\setminus \{0\}} \frac{x^*(i\omega I_n -JQ)^*(i\omega I_n -JQ)x}{x^*Q^2x}\\
%	& = (\lambda_{\text{min}}(R))^2 + \inf_{\omega \in \R} \lambda_{\text{min}}(Q^{-1}(i\omega I_n -JQ)^*(i\omega I_n -JQ)Q^{-1}),
%\end{align*}
%%
%where in the last equation, the infimum in the first part is obtained by the vector $\hat x = Q^{-1}x_0$, where $x_0$ is the eigenvector corresponding to the smallest eigenvalue of $R$. Next we provide a corollary to Theorem~\ref{thm:def}, which gives an upper bound for $r^{S_d}(J,R)$.
%
% \begin{theorem}\label{cor:def}
% 	Consider an asymptotically stable dHDAE system of the form~\eqref{dhdae1}. Let  $P^*RP=D$ be the spectral decomposition of $R$, where $P\in \C^{n,n}$ is unitary and $D=\text{diag}(d_1,d_2,\ldots,d_n)$ is such that $d_1\geq d_2\geq \ldots\geq d_n\geq 0$. Then
% 	%
% 	\begin{equation}\label{eq:defbound}
% 		d_n^2 + \inf_{\omega \in \R} \lambda_{\min}((i\omega E -J)^*(i\omega E -J)) \leq d_{im}^{\mathcal S_d}(J,R)^2\leq d_n^2 + \inf_{\omega\in\R} \frac{\|(i\omega E-J)\hat x\|^2}{\|\hat x\|^2},
% 	\end{equation}
% 	%
% 	where $\hat x=Pe_n$, with $e_n$ being the nth column of the $n\times n$ identity matrix.
% \end{theorem}

The situation is different for $E\geq 0$ with $E$ being singular. In contrast to the case $E>0$ we now can have loss of regularity of the pencil, an increase in the index, or the presence of a finite eigenvalue on the imaginary axis.
In the following, we first focus on moving finite eigenvalues to the imaginary axis, considering only structure-preserving perturbations in $J$ and $R$ while keeping $E$ fixed. The following result, provides an explicit characterization of this distance.
%
\begin{theorem}\label{thm:dist-imJR}
	Consider a robustly asymptotically stable dHDAE system of the form~\eqref{dhdae1}. Then
	%
	\begin{enumerate}
		\item the distance to the nearest dH pencil with purely imaginary eigenvalues $d_{im}^{\mathcal S_d}(J,R)$ is given by
		\begin{equation}\label{eq:dist-imp}
			(d_{im}^{\mathcal S_d}(J,R))^2 = \inf_{\omega \in \R} \inf_{x\in \C^{n}\setminus \{0\}} \left\{\left(\frac{x^*R^2x}{x^*Rx} \right)^2+\frac{x^*(i\omega E -J)^*(i\omega E -J)x}{x^*x} \right\},
		\end{equation}
		\item if $R>0$, the distance to the nearest dH pencil with purely imaginary eigenvalues $d_{im}^{\mathcal S_i}(J,R)$ is given by
		\begin{equation*}
			d_{im}^{\mathcal S_i}(J,R)=\inf_{\omega\in \R} \sigma_{\min}\left( \mat{c}R\\i\omega E - J \rix \right)
		\end{equation*}
		and, if $R\geq 0$ is singular, then
		\begin{equation*}
			d_{im}^{\mathcal S_i}(J,R)\geq \inf_{\omega\in \R} \sigma_{\min}\left( \mat{c}R\\i\omega E - J \rix \right).
		\end{equation*}
	\end{enumerate}
\end{theorem}
\begin{proof}
	By Definition~\ref{def:distances}, we have
	%
	\begin{align*}
		d_{im}^{\mathbb S}(J,R)=\inf \{ \nnrm{(\Delta_J,\Delta_R)}~:~(\Delta_J,\Delta_R)\in\mathbb S,~\Lambda(E,J+\Delta_J - (R+\Delta_R)) \cap i\R \neq \emptyset \}.
	\end{align*}
	%
	Since for $(\Delta_J,\Delta_R)\in \mathbb S\in \{\mathcal S_d(J,R),\mathcal S_i(J,R)\}$, the perturbed dH pair $(E,J+\Delta_J - (R+\Delta_R))$ still has the dH structure, by using~\cite[Lemma 4.1]{MehMS16}, we obtain
	%
	\begin{align}%\label{eq:def1}
		d_{im}^{\mathbb S}(J,R)
		&=\inf \{ \nnrm{(\Delta_J,\Delta_R)}:(\Delta_J,\Delta_R)\in\mathbb S,~(R+\Delta_R)x=0 \text{ for some eigenvector $x$ of } (E,J+\Delta_J) \}\nonumber\\
		&=\inf\{ \nnrm{(\Delta_J,\Delta_R)}:(\Delta_J,\Delta_R)\in\mathbb S,~ (R+\Delta_R)x=0 \text{ for some $x\in \C^n\setminus\{0\}$ satisfying } \nonumber\\ &\hspace{5cm}  (J+\Delta_J)x= i\omega E x, \omega\in \R\}\nonumber\\
		&= \inf_{\omega\in \R}\inf_{x\in\C^{n}\setminus \{0\}} \Big(\inf \{ \nnrm{(\Delta_J,\Delta_R)}:(\Delta_J,\Delta_R)\in\mathbb S,~ (R+\Delta_R)x=0, (J+\Delta_J)x= i\omega E x\} \Big)\nonumber\\
		&=\inf_{\omega \in \R } \vartheta^{\mathbb S}_{\omega} \label{eq:def11},
	\end{align}
	%
	%    where
	%the equation~\eqref{eq:def11} separated the variable $x$, $%\omega$ and $(\Delta_J,\Delta_R)$ into sublevel optimizations.
	where, for a given scalar $\omega$, we have
	%
	\begin{align}\label{eq:def12}
		\vartheta^{\mathbb S}_{\omega}:= \inf_{x\in\C^{n}\setminus \{0\}} \Big(\inf \{ \nnrm{(\Delta_J,\Delta_R)}:(\Delta_J,\Delta_R)\in\mathbb S,~ (R+\Delta_R)x=0,(J+\Delta_J)x= i\omega E x\} \Big).
	\end{align}
	%
	Let $\mathbb S=\mathcal S_d(J,R)$. Then, for the inner optimization in~\eqref{eq:def12}, we show that the minimal value can be expressed as the sum of two generalized Rayleigh quotients, which depend on the variables $x$ and $\omega$. For this, we first solve two mapping problems.

    By Theorem~\ref{map:herm} and~\cite[Theorem~2.3]{MehMS16}, there exist a skew-Hermitian matrix $\Delta_J$ and a Hermitian negative semidefinite matrix $\Delta_R\le 0$ satisfying the constraints in~\eqref{eq:def12} if and only if
    \[
    x^*(i\omega E - J)x \in i\mathbb{R}
    \quad \text{and} \quad
    -x^*Rx<0.
    \]
    These conditions are automatically satisfied due to the structural properties of the matrices $E,~J$, and $R$.
    %
    Among all such admissible mappings, the minimal norms are given by
    %
    \begin{equation}\label{norm:jdef}
		\|\Delta_J\|=\frac{\|(i\omega E - J)x\|}{\|x\|}, \quad \text{and} \quad
		\|\Delta_R\|=\frac{\|Rx\|^2}{x^*Rx},
	\end{equation}
	%
    and are attained by $\Delta_J=-i \hat H_{(x, i(i\omega E-J)x)}$, where $\hat H$ is defined in~\eqref{def:hatH}, and $\Delta_R=\frac{-1}{x^*Rx}(Rx)(Rx)^*$ (if $Rx\neq 0$). If $Rx=0$ (this case may arise if $R$ is singular), then again from {Theorem~\ref{map:def}}, $\Delta_R=0$ is the solution that satisfies~\eqref{norm:jdef} by defining the fraction $\frac{0}{0}$ to have value $0$. Note that while solving the mapping constraint $\Delta_Rx=-Rx$, we have not considered the additional condition $R+\Delta_R\geq 0$. Using the fact that $\mathcal S_d(J,R)\subseteq \{(\Delta_J,\Delta_R): \Delta_J^*=-\Delta_J, \Delta_R^*=\Delta_R\leq 0\}$, and using~\eqref{norm:jdef} in~\eqref{eq:def12}, we obtain
	%
	%	Similarly for $\Delta_R$, on using Lemma~\ref{map:def}, there exists $\Delta_R<0$ satisfying $\Delta_RQx=-RQx$ if and only if $-x^*QRQx<0$, which hold always because $R$ is Hermitian positive definite matrix. The minimal norm mapping $\Delta_R$ is given by
	%	%
	%	\begin{equation}\label{norm:rdef}
		%		\|\Delta_R\|=\frac{\|RQx\|^2}{x^*QRQx}
		%	\end{equation}
	%	%
	%	which is attained by $\Delta_R=\frac{1}{x^*QRQx}(RQx)(RQx)^*$.
	%	%
	%	Using~\eqref{norm:jdef} and~\eqref{norm:rdef} and the fact that $S_d(J,R)\subseteq \{(\Delta_J,\Delta_R): \Delta_J^*=-\Delta_J, \Delta_R^*=\Delta_R<0\}$, we obtain
	%	
	\begin{equation}\label{eq:def2}
		(\vartheta^{\mathcal S_d}_{\omega})^2\geq \inf_{x\in\C^{n}\setminus\{0\}} \left\{ \frac{\|Rx\|^4}{(x^*Rx)^2}+\frac{\|(i\omega E - J)x\|^2}{\|x\|^2} \right\}.
	\end{equation}
	%
	% Next, we show that equality holds in~\eqref{eq:def2}. Note that the vector $x$ can be scaled to have unit norm without affecting the objective function in~\eqref{eq:def2}. Using the compactness of the set of unit norm vectors, the infimum in~\eqref{eq:def2} is actually a minimum and is attained for some $x=\hat x$. For a given $\omega$, set
	% %
	% \begin{equation}
	% 	\hat{\Delta}_R=\begin{cases}
	% 		-\frac{(R\hat x)(R\hat x)^*}{\hat x^*R\hat x} & \text{ if } R\hat{x} \neq 0,\\
	% 		\qquad 0 & \text{ if } R\hat{x} = 0,
	% 	\end{cases} \quad \quad \hat \Delta_J = \hat H_{(\hat x,(i\omega E - J)\hat x)},
	% \end{equation}
	% %
	% where $\hat H$ is defined in~\eqref{def:hatH}. By {\cite[Lemma 4.1]{MehMS16}}, it then follows that for the optimal perturbation $\hat \Delta_R$ we have that $R+\hat \Delta_R\geq0$ so that equality holds in~\eqref{eq:def2}, as then $(\hat \Delta_J,\hat \Delta_R)\in \mathcal S_d(J,R)$.
    %
    %
    Next we consider $\mathbb S=\mathcal S_i(J,R)$. By Theorem~\ref{map:herm}, there exist a skew-Hermitian matrix $\Delta_J$ and a Hermitian matrix $\Delta_R$ satisfying the constraints in~\eqref{eq:def12} if and only if
    \[
    x^* (i\omega E - J)x \in i\mathbb{R}
    \quad \text{and} \quad
    x^* R x \in \mathbb{R},
    \]
    which again holds trivially due to the structure of $E,~J$ and $R$.
    %
    The corresponding minimal norms are
    \[
    \|\Delta_J\| = \frac{\|(i\omega E - J)) x\|}{\|x\|},
    \qquad
    \|\Delta_R\| = \frac{\|R x\|}{\|x\|},
    \]
    and are attained by
    \[
    \Delta_J = -i\,\hat H_{(x, i(i\omega E - J) x)},
    \qquad
    \Delta_R = -\hat H_{(x,Rx)},
    \]
    where $\hat H$ is defined in~\eqref{def:hatH}. Note that $(R+\Delta_R)x=0$, hence, by applying~\cite[Lemma 4.4]{MehMS16} we have $R+\Delta_R\geq 0$, implying that $(\Delta_J,\Delta_R)\in \mathcal S_i(J,R)$. Using the above arguments in~\eqref{eq:def12}, we obtain
    %
    \begin{equation*}
		\vartheta^{\mathcal S_i}_{\omega} = \inf_{x\in\C^{n}\setminus\{0\}} \left\{ \frac{\sqrt{\|Rx\|^2+\|(i\omega E - J)x\|^2}}{\|x\|}\right\} = \sigma_{\min}\left(\mat{c}R\\i\omega E - J \rix\right).
	\end{equation*}
    %
    This completes the proof.
	%This follows directly by {\cite[Lemma 4.1]{MehMS16}} %on %using Lemma~\ref{lemdh4}. This completes the proof.
	%
	% {\color{red} Proof can be shortened or removed by citing the previous paper. Proof of the second part $d_{im}^{S_i}$ also follows on the similar lines.
 %    This case can be transforemed to the standard DH system in the form $(J-R)Q$, where $Q=E^{-1}$.}
\end{proof}
%
 In the following, we state a result for the structured distances to singularity $d_{sing}^{\mathcal S_d}(J,R)$ and $d_{sing}^{\mathcal S_i}(J,R)$ keeping $E$ fixed. We skip its proof as the distance $d_{sing}^{\mathcal S_d}(J,R)$ was obtained in~\cite[Table A.1]{PraS22} and is restated here adopting the convention that $\frac{0}{0}=0$. On the other hand, the proof for $d_{sing}^{\mathcal S_i}(J,R)$ follows by arguments similar to the proof of $d_{sing}^{\mathcal S_i}(E,J,R)$ in Theorem~\ref{thm:dist_singEJR_i}.


\begin{theorem}\label{thm:dist-singJR}
	Consider a robustly asymptotically stable dHDAE system of the form~\eqref{dhdae1}. Then, the distance to singularity $d_{sing}^{\mathbb S}(J,R)$ with respect to perturbations from the set $\mathbb S\in \{\mathcal S_d(J,R),\mathcal S_i(J,R)\}$ is $\infty$, if $E>0$. If $E\geq 0$ is singular, then let the columns of the matrix $N$ form a basis of the null space of $E$. Then
	%
    \begin{enumerate}
        \item the distance to singularity $d_{sing}^{\mathcal S_d}(J,R)$ with respect to perturbations from the set $\mathcal S_d$ is
        %
    	\begin{equation}\label{eq:dist-sing}
    		(d_{sing}^{\mathcal S_d}(J,R))^2=
    			\inf_{x\in \C^{r}\setminus \{0\}} \left\{\left(\frac{x^*N^*R^2Nx}{x^*N^*RNx} \right)^2+\frac{x^*N^*J^*JNx}{x^*x} \right\},
    	\end{equation}
        %
        \item if $R>0$, then the distance to singularity $d_{sing}^{\mathcal S_i}(J,R)$ with respect to perturbations from the set $\mathcal S_i(J,R)$ is
        %
        \begin{equation}
            d_{sing}^{\mathcal S_i}(J,R) = \sqrt{\lambda_{\min}(N^*(R^2-J^2)N)},
        \end{equation}
        %
        and, if $R\ge 0$ is singular, then
        %
        \begin{equation}
            d_{sing}^{\mathcal S_i}(J,R) \geq \sqrt{\lambda_{\min}(N^*(R^2-J^2)N)}.
        \end{equation}
    \end{enumerate}
\end{theorem}
% \begin{proof}
% 	By Definition~\ref{def:distances},
% 	%
% 	\begin{align*}
% 		&d_{sing}^{\mathbb S}(J,R)\\
% 		&=\inf \left\{ \nnrm{(\Delta_J,\Delta_R)}~:~(\Delta_J,\Delta_R)\in \mathbb S,~ \ker(E) \cap \ker(J + \Delta_J) \cap \ker(R+\Delta_R) \neq \{0\} \right\}.\\
% 		&=\inf \left\{ \nnrm{(\Delta_J,\Delta_R)}~:~(\Delta_J,\Delta_R)\in\mathbb S,~x\in \C^{n}\setminus\{0\},~ Ex=0, (J + \Delta_J)x=0, (R+\Delta_R)x=0 \right\}.
% 	\end{align*}
% 	%
% 	Clearly, if $\ker (E)=\{0\}$, then $d_{sing}^{\mathbb S}(J,R)=\infty$. We therefore restrict our attention to the nontrivial case ker$(E)\neq \{0\}$. Fix $x\in \C^{n}\setminus \{0\}$ and consider the problem of determining a minimum-norm perturbation $(\Delta_J,\Delta_R)\in \mathbb S$ such that
% 	%
% 	\begin{equation*}
% 		Ex=0,~ \Delta_Rx=-Rx, \text{ and }~ \Delta_Jx=-Jx.
% 	\end{equation*}
% 	%
% 	Let the columns of $U\in \mathbb C^{n,r}$ form an orthonormal basis of $\ker (E)$. Then $Ex=0$ if and only if $x=Uy$ for some $y\in \mathbb C^{r}\setminus \{0\}$. With this parametrisation, the above problem reduces to
% 	%
% 	\begin{equation}\label{eq:dist-sing1}
% 		d_{sing}^{\mathbb S}(J,R)=\inf \left\{ \nnrm{(\Delta_J,\Delta_R)}~:~(\Delta_J,\Delta_R)\in\mathbb  S,~y\in \C^{r}\setminus\{0\},~ (J + \Delta_J)Uy=0, (R+\Delta_R)Uy=0 \right\}.
% 	\end{equation}
% 	%
%     When $\mathbb S=\mathcal S_i(J,R)$, then by Theorem~\ref{map:herm}, there exists a skew-Hermitian perturbation $\Delta_J$ satisfying $\Delta_JUy=- JUy$ and Hermitian perturbation $\Delta_R$ satisfying $\Delta_RUy=-RUy$ if and only if $y^*U^*JUy\in i\R$ and $y^*U^*RUy\in \R$, respectively, which holds trivially due to the skew-Hermitian Hermitian structure of $J$ and $R$. Among all such perturbations, the minimal norm is given by
% 	%
% 	\begin{equation}\label{norm:jdef-sing}
% 		\|\Delta_J\|=\frac{\|JUy|}{\|Uy\|}, \quad \text{and} \quad \|\Delta_R\|=\frac{\|RUy|}{\|Uy\|}
% 	\end{equation}
% 	%
% 	and this minimum is attained for $\Delta_J=i \hat H_{(Uy, -iJUy)}, \Delta_R=\hat H_{(Uy,-RUy)}$, where $\hat H$ is defined in \eqref{def:hatH}.
	
%     % If $RUy=0$ (which may occur when $R$ is singular), then again by \cite[Theorem~2.3]{MehMS16}, the choice $\Delta_R=0$ satisfies the constraint $\Delta_RUy=-RUy$. In this case, \eqref{norm:rdef-sing} remains valid by adopting the convention 0/0=00/0=0.
%     We note that, while considering the mapping constraint $\Delta_RUy=-RUy$, we have not explicitly imposed the additional requirement $R+\Delta_R\geq 0$. But using the fact that, $\mathcal S_i(J,R)\subseteq \{(\Delta_J,\Delta_R): \Delta_J^*=-\Delta_J, \Delta_R^*=\Delta_R\leq 0\}$, and combining \eqref{norm:jdef-sing} and \eqref{norm:rdef-sing} in \eqref{eq:dist-sing1}, we obtain
% 	%
% 	\begin{equation}\label{eq:dist-sing2}
% 		(d_{sing}^{\mathcal S_i}(J,R))^2\geq \inf_{y\in\C^{r}\setminus\{0\}} \left\{ \frac{y^*U^*(R^*R+J^*J)Uy}{y^*y}\right\}=\lambda_{\min}(U^*(R^2-J^2)U).
% 	\end{equation}
% 	%
%     Suppose $R>0$, then if we can show that equality holds in~\eqref{eq:dist-sing2}, then the proof is finished. Note that the vector $y$ can be scaled to have unit norm without affecting the objective function in~\eqref{eq:dist-sing2}. Using the compactness of the set of unit norm vectors, the infimum in~\eqref{eq:dist-sing2} is actually a minimum and is attained for some $y=\hat y$. Set
% 	%
% 	% \begin{equation}
% 	% 	\hat{\Delta}_R=\begin{cases}
% 	% 		-\frac{(RU\hat y)(RU\hat y)^*}{\hat y^*U^*RU\hat y} & \text{ if } RU\hat{y} \neq 0,\\
% 	% 		\qquad 0 & \text{ if } RU\hat{y} = 0,
% 	% 	\end{cases} \quad \quad \hat \Delta_J = \hat H_{(U\hat y,- JU\hat y)},
% 	% \end{equation}
% 	% %
% 	% where $\hat H$ is defined in~\eqref{def:hatH}. By {\cite[Lemma 4.1]{MehMS16}}, it then follows that for the optimal perturbation $\hat \Delta_R$ we have that $R+\hat \Delta_R\geq0$ so that equality holds in~\eqref{eq:dist-sing2}, as then $(\hat \Delta_J,\hat \Delta_R)\in S_d(J,R)$.
% \end{proof}


 In the following theorem, we study the distance to higher index, allowing structure-preserving perturbations only in $J$ and $R$ while keeping $E$ fixed.

\begin{theorem}\label{thm:dist-hiJR}
    Consider an asymptotically stable dHDAE system of the form~\eqref{dhdae1}. Then
    %
    \begin{enumerate}
        \item the distance to higher index $d_{hi}^{\mathcal S_d}(J,R)$ with respect to perturbations from the set $\mathcal S_{d}(J,R)$ is given by
        %
        \begin{equation}
            (d_{hi}^{\mathcal S_d}(J,R))^2 = \inf_{x\neq 0} \left\{\frac{\|(N^*)^{\dagger} N^* J N x\|^2}{\|N x\|^2} + \frac{\|(N^*)^{\dagger} N^* R N x\|^4} {\big(x^* N^* R N x\big)^2} \right\},
        \end{equation}
        %
        where $N$ denotes a unitary  matrix whose columns form a basis for $\ker(E)$,
        \item the distance to higher index $d_{hi}^{\mathcal S_i}(J,R)$ with respect to perturbations from the set $\mathcal S_{i}(J,R)$ is given by
        %
        \begin{equation}
            (d_{hi}^{\mathcal S_i}(J,R))^2 = \inf_{x\neq 0} \left\{ \frac{\|(N^*)^{\dagger} N^* J N x\|^2}{\|N x\|^2} + \frac{\|(N^*)^{\dagger} N^* R N x\|^2} {\|N x\|^2} \right\},
        \end{equation}
        %
        where $N$ denotes a unitary matrix whose columns form a basis for $\ker(E)$.
    \end{enumerate}
\end{theorem}
\begin{proof}
    From Definition~\eqref{def:distances}, the distance to higher index under structured perturbations from the set $\mathbb S\in \{\mathcal S_d(J,R),\mathcal S_i(J,R)\}$ is given by
    \[
    d_{hi}^{\mathbb S}(J,R) = \inf\Big\{\nnrm{(\Delta_J,\Delta_R)} ~:~ (\Delta_J,\Delta_R)\in \mathbb S,~ \text{ index of } (E,J+\Delta_J-R-\Delta_R) > 1 \Big\}.
    \]
    Equivalently,
    %
    \begin{equation*}
        d_{hi}^{\mathbb S}(J,R) = \inf\Big\{ \nnrm{(\Delta_J,\Delta_R)} ~:~ (\Delta_J,\Delta_R)\in \mathbb S,~ \text{rank}\!\big(N(E)^* (J+\Delta_J-R-\Delta_R) N(E) \big) < n \Big\},
    \end{equation*}
    %
    where $N$ denotes an orthonormal basis matrix for the null space of $E$.
%
    Note that $d_{hi}^{\mathbb S}(J,R)=\mu_k^{\mathbb S}$, where $\mu_k^{\mathbb S}$ is defined by~\eqref{hu_hi} for $N_k=N$. Thus the proof follows by using the expressions of $\mu_k^{\mathcal S_d}$ and $\mu_k^{\mathcal S_i}$ from~\eqref{Sd:high11} and~\eqref{Sd:high22}, respectively.\end{proof}
    % Since $J+\Delta_J$ is skew-Hermitian and $R+\Delta_R$ is Hermitian positive semidefinite, the condition
    % \[
    % (J+\Delta_J - R - \Delta_R)y = 0
    % \]
    % holds if and only if
    % \[
    % (J+\Delta_J)y = 0 \quad \text{and} \quad (R+\Delta_R)y = 0.
    % \]
    % Therefore,
    % %
    % \begin{align*}
    %     d_{hi}^{\mathbb S}(J,R) = \inf\Big\{\nnrm{(\Delta_J,\Delta_R)} : (\Delta_j,\Delta_R)\in \mathbb S,\exists\, x\neq 0 \text{ s.t. } N^* (J+\Delta_J) N x = 0,~ N^* (R+\Delta_R) N x = 0 \Big\}.
    % \end{align*}
    % %
    % Note from~\cite[Lemma 2.8]{MehMS16} that for a matrix $B$, the equation $B\Delta x = y$ holds if and only if $\Delta x = B^{\dagger}y$ and $BB^{\dagger}y = y$, where $B^{\dagger}$ denotes the Moore-Penrose pseudoinverse. Applying this to the above constraints, the relations
    % \[
    % N^* \Delta_J N x = - N^* J N x, \qquad N^* \Delta_R N x = - N^* R N x
    % \]
    % can be equivalently written as
    % \[
    % \Delta_J N x = - (N^*)^{\dagger} N^* J N x, \qquad \Delta_R N x = - (N^*)^{\dagger} N^* R N x,
    % \]
    % since the consistency conditions are trivially satisfied.

    % Consequently, we obtain
    % %
    % \begin{align}\label{Sd:highJR1}
    %     d_{hi}^{\mathbb S}(J,R) = & \inf\Big\{ \nnrm{(\Delta_J,\Delta_R)} : (\Delta_J,\Delta_R)\in \mathbb S,~\exists\, x\neq 0, \; \Delta_J N x = - (N^*)^{\dagger} N^* J N x,\nonumber \\ & \hspace{4cm}\Delta_R N x = - (N^*)^{\dagger} N^* R N x \Big\}.
    % \end{align}
    % %
    % We now solve the associated minimal-norm mapping problems.
    % First, let $\mathbb S=\mathcal S_d(J,R)$. By Theorem~\ref{map:herm} and~\cite[Theorem~2.3]{MehMS16}, there exist a skew-Hermitian matrix $\Delta_J$ and a Hermitian negative semidefinite matrix $\Delta_R\le 0$ satisfying the constraints in~\eqref{Sd:highJR1} if and only if
    % \[
    % x^* N^* J N x \in i\mathbb{R}
    % \quad \text{and} \quad
    % x^* N^* R N x > 0.
    % \]
    % These conditions are automatically satisfied due to the structural properties of the matrices $J$ and $R$.
    % %
    % Among all such admissible mappings, the minimal norms are given by
    % \[
    % \|\Delta_J\| = \frac{\|(N^*)^{\dagger} N^* J N x\|}{\|N x\|}, \qquad
    % \|\Delta_R\| = \frac{\|(N^*)^{\dagger} N^* R N x\|^2}{x^* N^* R N x},
    % \]
    % and are attained by
    % \[
    % \Delta_J = i\,\hat H_{(N x,i(N^*)^{\dagger} N^* J N x)}, \qquad
    % \Delta_R = -\frac{\big((N^*)^{\dagger} N^* R N x\big) \big((N^*)^{\dagger} N^* R N x\big)^*}{x^* N^* R N x},
    % \]
    % where $\hat H$ is defined in~\eqref{def:hatH}.
    % %
    % If $R\ge 0$ is singular and $x^* N^* R N x = 0$, then necessarily $(N^*)^{\dagger} N^* R N x = 0$. In this case, adopting the convention $0/0 = 0$, the choice $\Delta_R = 0$ satisfies both the mapping constraint in~\eqref{Sd:highJR1} and the minimal-norm condition. Moreover, the resulting perturbation satisfies $R + \Delta_R \ge 0$; see~\cite[Lemma~4.1]{MehMS16}.
    % %
    % Combining these observations, we obtain
    % \[
    % (d_{hi}^{\mathcal S_d}(J,R))^2 = \inf_{x \neq 0} \left\{
    % \frac{\|(N^*)^{\dagger} N^* J N x\|^2}{\|N x\|^2} +
    % \frac{\|(N^*)^{\dagger} N^* R N x\|^4}{\big(x^* N^* R N x\big)^2} \right\}.
    % \]
    % Next, consider $\mathbb S=\mathcal S_i(J,R)$. By Theorem~\ref{map:herm}, there exist a skew-Hermitian matrix $\Delta_J$ and a Hermitian matrix $\Delta_R$ satisfying the constraints in~\eqref{Sd:highJR1} if and only if
    % \[
    % x^* N^* J N x \in i\mathbb{R}
    % \quad \text{and} \quad
    % x^* N^* R N x \in \mathbb{R},
    % \]
    % which again holds trivially due to the structure of $J$ and $R$.
    % %
    % The corresponding minimal norms are
    % \[
    % \|\Delta_J\| = \frac{\|(N^*)^{\dagger} N^* J N x\|}{\|N x\|},
    % \qquad
    % \|\Delta_R\| = \frac{\|(N^*)^{\dagger} N^* R N x\|}{\|N x\|},
    % \]
    % and are attained by
    % \[
    % \Delta_J = i\,\hat H_{(N x,\, i (N^*)^{\dagger} N^* J N x)},
    % \qquad
    % \Delta_R = -\hat H_{(N x,\,(N^*)^{\dagger} N^* R N x)},
    % \]
    % where $\hat H$ is defined in~\eqref{def:hatH}. Note that $(R+\Delta_R)Nx=0$, hence on applying~\cite[Lemma 4.4]{MehMS16} we have $R+\Delta_R\geq 0$.
    % %
    % Applying these arguments to~\eqref{Sd:highJR1} yields
    % \[
    % (d_{hi}^{\mathcal S_i}(J,R))^2 \ge \inf_{x \neq 0} \left\{
    % \lambda_{n-k+1}(E)^2 +
    % \frac{\|(N^*)^{\dagger} N^* J N x\|^2}{\|N x\|^2} +
    % \frac{\|(N^*)^{\dagger} N^* R N x\|^2}{\|N x\|^2} \right\}.
    % \]
\begin{table}[h!]\label{tablesummary}
\centering
\renewcommand{\arraystretch}{1.5}
\begin{tabular}{lccc}
\hline
\textbf{$\mathbb  S$}
& $\mathbf{d_{im}^{\mathbb S}}$
& $\mathbf{d_{sing}^{\mathbb S}}$
& $\mathbf{d_{hi}^{\mathbb S}}$ \\
\hline
$\mathcal S_d(E,J,R)$
& Lower bound: Theorem~\ref{thm:dist-imagEJR}
& Theorem~\ref{thm:dist-singEJR-d}
& Upper bound: Theorem~\ref{thm:dist-hiEJR} \\
%
$\mathcal S_i(E,J,R)$
& Lower bound: Theorem~\ref{thm:dist-imagEJR}
& Theorem~\ref{thm:dist_singEJR_i}
& Upper bound: Theorem~\ref{thm:dist-hiEJR} \\
%
$\mathcal S_d(J,R)$
& Theorem~\ref{thm:dist-imJR}
& Theorem~\ref{thm:dist-singJR}
& Theorem~\ref{thm:dist-hiJR} \\

$\mathcal S_i(J,R)$
& Theorem~\ref{thm:dist-imJR}
& Theorem~\ref{thm:dist-singJR}
& Theorem~\ref{thm:dist-hiJR} \\
\hline
\end{tabular}
\caption{Summary of structured distances and corresponding results}
\end{table}
%Put table for all cases.
%\begin{itemize}
%\item $E>0, \Delta_E=0$, $\Delta_R>0$, $\Delta_R$ indefinite.
%\item $E>0, \Delta_E\neq0$, $E+\Delta_E>0$, $\Delta_R>0$, $\Delta_R$ indefinite.
%\item $E>0, \Delta_E\neq0$, $E+\Delta_E\geq 0$, $\Delta_R>0$, $\Delta_R$ indefinite. All Three cases a), b), c)
%\item $E\geq 0, \Delta_E=0$, $E+\Delta_E\geq 0$, $\Delta_R>0$, $\Delta_R$ indefinite. All Three cases a), b), c)
%\item $E\geq 0, \Delta_E\neq0$, $E+\Delta_E\geq 0$, $\Delta_R>0$, $\Delta_R$ indefinite. All Three cases a), b), c)
%\end{itemize}
%{\color{red} Put detailed discussion here, discuss on monday} upper and lower bounds, mapping theorems, ....


In this section we have presented
bounds and analytic expressions for  the different structured distances. They are summarized in Table~\ref{tablesummary}. In the next section we illustrate these results via  numerical examples.

\section{Numerical examples}\label{sec:numerical}
In this section, we present numerical experiments to illustrate the theoretical results developed in the previous sections. In particular, we demonstrate the computation of the structured distances to systems with purely imaginary  eigenvalues, that are singular, or have higher index for dHDAEs. All computations are carried out in \textsc{Matlab} R2024b. %

In the following, we briefly describe the numerical procedures that are used to compute the various structured stability radii. For several of the distances derived in the previous sections, namely
$d_{sing}^{\mathcal S_d}(E,J,R)$, $d_{sing}^{\mathcal S_d}(J,R)$, $d_{hi}^{\mathcal S_d}(E,J,R)$, $d_{hi}^{\mathcal S_i}(E,J,R),$ $d_{hi}^{\mathcal S_d}(J,R)$, and $d_{hi}^{\mathcal S_i}(J,R)$, the characterization reduces to an optimization problem involving the sum of two generalized Rayleigh quotients. More precisely, these quantities can be written in the form
\[
\inf_{\omega \in \mathbb{R}} \; \inf_{x \neq 0} \left\{ \frac{x^* H_1 x}{x^* x} + \left(\frac{x^* H_2 x}{x^* H_3 x}\right)^2 \right\},
\]
where $H_1, H_2,$ and $H_3$ are Hermitian positive semidefinite matrices (depending on~\(\omega\)).
%
To evaluate such expressions numerically, we employ a two-level optimization strategy. For the inner minimization with respect to $x$, we use the nonlinear eigenvalue problem (NEPv) characterization and solve the resulting NEPv via level-shifted self-consistent field (SCF) iterations; see~\cite{LuPSB25} for details on this approach. The outer minimization with respect to~\(\omega\) is then carried out using the function \texttt{fminsearch} in \textsc{Matlab}.

In the remaining cases, the computations are simpler. The distances
$d_{sing}^{\mathcal S_i}(E,J,R)$, and $d_{sing}^{\mathcal S_i}(J,R)$ reduce to standard eigenvalue problems and are computed directly using built-in eigenvalue solvers. The quantities $d_{im}^{\mathcal S_i}(E,J,R)$ and $d_{im}^{\mathcal S_i}(J,R)$ lead to a parameter-dependent eigenvalue problems, which we handle by combining eigenvalue computations with an outer minimization over the parameter using \texttt{fminsearch}.

Since the computation of $d_{im}^{\mathcal S_d}(E,J,R)$ involves a constrained optimization problem, we use \textsc{Matlab}’s \texttt{fmincon} function. As \texttt{fmincon} may not always provide highly accurate estimates for the solution of thiis type of problems, a possible direction for future work is to develop more reliable numerical methods by utilizing the special structure of the optimization.

The codes and data of the examples presented below are available at \url{https://gitlab.mpi-magdeburg.mpg.de/prajapati/dhdae-stability-radii.git}.

%The only case for which we do not present numerical examples  is $d_{im}^{\mathcal S_d}(E,J,R),$ for which an efficient and reliable computational procedure is currently not available. We  leave this topic for future work.

% In many cases, the characterizations of the various structured stability radii obtained in the previous sections reduce to optimization problems involving the sum of two generalized Rayleigh quotients; see, for instance, Theorems~\ref{thm:dist-hiEJR}, \ref{thm:dist-imJR},  \ref{thm:dist-singJR}, and \ref{thm:dist-hiJR}. In each of these results, the corresponding reformulation leads to an equivalent optimization problem of the form
% \[
% \inf_{\omega\in \R}\inf_{x \neq 0} \frac{x^*H_1x}{x^*x} + \left(\frac{x^*H_2x}{x^*H_3x}\right)^2,
% \]
% where $H_1, H_2$ and $H_3$ are Hermitian positive semidefinite matrices depending on the variable $\omega$.

% To compute such quantities numerically, we adopt a two-level optimization strategy. For the inner optimization, we exploit the nonlinear eigenvalue problem (NEPv) characterization of the stationary points and solve the associated NEPv using level-shifted self-consistent field (SCF) iterations; see~\cite{LuPSB25} for details on this approach. The outer optimization is then solved using the function \texttt{fminsearch} in \textsc{Matlab}.

%BeaMXZ18_example13_8x8
\begin{example}{\rm
    Consider Example~13 from~\cite{BeaMXZ18}, which arises from the finite element modeling of the acoustic field in the interior of a car. After simplifications, the associated differential–algebraic system is of the form
    \[
    M \ddot{p} + D \dot{p} + K p = B_1 u,
    \]
    where $p$ denotes the vector of coefficients corresponding to the acoustic pressure in the air and the structural displacements. The term $B_1 u$ represents an external force. The matrix $M$ is a symmetric positive semidefinite mass matrix, $D$ is a symmetric positive semidefinite damping matrix, and $K$ is a symmetric positive definite stiffness matrix. The resulting first-order formulation yields the state equation of a pHDAE system, $E \dot{x} = (J-R) x + B u$, where the state vector is $x = \mat{c}\dot{p}\\ p \rix $ and the matrices are given by
    %
    \[
    E = \mat{cc} M & 0 \\ 0 & K \rix,
    \quad
    J = \mat{cc} 0 & -K \\ K & 0 \rix,
    \quad
    R = \mat{cc} D & 0 \\ 0 & 0 \rix, \text{ and}
    \quad
    B = \mat{c} B_1 \\ 0 \rix.
    \]
    %
    For the numerical experiment, we randomly generated the matrices $M, K$, and $D\in \C^{4,4}$ in \textsc{Matlab} and enforce symmetry (and positive definiteness where required). We construct $M \in \mathbb{C}^{4\times 4} $ with rank $2$, so that the matrix $E$ becomes singular and, consequently, the system possesses two infinite eigenvalues. For the generated example,  using the criteria from Section~\ref{sec:distalg},  we verified that the resulting dHDAE system is robustly asymptotically stable. We then computed all the structured distances introduced in Section~\ref{sec:partial}; the resulting values are summarized in Table~\ref{tab:example1} below.
    %
    \begin{table}[h!]
        \centering
        \renewcommand{\arraystretch}{1.5}
        \begin{tabular}{cccc}
            \hline
            \textbf{$\mathbb  S$}
            & $\mathbf{d_{im}^{\mathbb S}(J,R)}$
            & $\mathbf{d_{sing}^{\mathbb S}(J,R)}$
            & $\mathbf{d_{hi}^{\mathbb S}(J,R)}$ \\

            & Theorem~\ref{thm:dist-imJR}
            & Theorem~\ref{thm:dist-singJR}
            & Theorem~\ref{thm:dist-hiJR}\\
            \hline
            $\mathcal S_d(J,R)$
            & 1.7631
            & 5.3705
            & 1.6546         \\

            $\mathcal S_i(J,R)$
            & 1.4351
            & 4.7242
            & 1.6546 \\
            \hline
        \end{tabular}
        \caption{Various structured distances.}
        \label{tab:example1}
    \end{table}
    %
    Recall that the distance to the boundary of the set of robustly asymptotically stable systems is defined as the minimum of the three distances listed in Table~\ref{tab:example1}. Hence, $d_{inst}^{\mathcal S_d}(J,R)=1.6546$ and $d_{inst}^{\mathcal S_i}(J,R)=1.4351$.
    %This quantity represents the smallest structured perturbation in $J$ and $R$ that destroys robust stability of the system.

    Interestingly, the mechanism through which the boundary of the robustly asymptotically stable region is first reached depends on the chosen perturbation set. When perturbations are restricted to the structure-preserving set $\mathcal S_d(J, R)$, the  nearest  system that is not robustly asymptotically stable is obtained due to an increase in the  index of the system, while the spectral abscissa being negative and the regularity are preserved. In contrast, for the structure-preserving set $\mathcal S_i(J, R)$, the minimal distance for a loss of robust asymptotic stability occurs due to an eigenvalue reaching the imaginary axis.
    %
    \begin{figure}[!h]
        \centering
        \includegraphics[width=0.8\textwidth]{fig_example1.png}
        \caption{Eigenvalue perturbation curves of dHDAE system with respect to structured complex perturbations from the set $\mathcal S_d(J,R)$ (left) and $\mathcal S_i(J,R)$ (right).}
        \label{fig:example1}
    \end{figure}
    %
    To illustrate the corresponding optimal perturbations $(\Delta_J,\Delta_R)$, we plot the eigenvalues of the pencil
    \[
    zE - (J - R)
    \]
    under the homotopy
    \[
    zE - (J + t\Delta_J - R - t\Delta_R),
    \qquad t \in [0,1],
    \]
    %
where $\Delta_J$ and $\Delta_R$ are the perturbations that give the minimal distance.
Figure~\ref{fig:example1} shows the eigenvalue trajectories when the optimal perturbations are taken from $\mathcal S_d(J,R)$~(left) and $\mathcal S_i(J,R)$~(right). The original eigenvalues are marked by stars surrounded by circles, and as $t$ increases from $0$ to $1$, one observes that the loss of robustness occurs due to an increase in the index of the system (left) and due to an eigenvalue on the imaginary axis (right). It turns out that the perturbations that lead to a loss of robust asymptotic stability have the form $\Delta_J=0 $ and $\Delta_R=\diag(\Delta,0)$, so that the zero block structure of the original system is preserved. In the left plot of Figure~\ref{fig:example1}, one observes that a pair of finite eigenvalues approaches the real axis, coalesces, and then splits: one branch moves toward infinity while the other moves toward zero. The migration of a finite eigenvalue to infinity signals  an increase in the multiplicity of the infinite eigenvalue. This behavior is characteristic for a transition to a higher index. This also explains why the eigenvalue plot corresponding to $\mathcal S_d(J,R)$ exhibits symmetry about the real axis.

When perturbations are taken from the larger set $\mathcal S_i(J,R)$, the minimal perturbation do not preserve the block structure of the original matrices. In this case, the eigenvalue nearest to the imaginary axis moves directly toward the imaginary axis and reaches it without preserving the symmetry.

For the same constructed example, we additionally computed all three structured distances with respect to perturbations from the full structure-preserving sets $\mathcal S_d(E,J,R) \text{ and }\mathcal S_i(E,J,R)$, where all three matrices $E$, $J$, and $R$ are allowed to vary. The corresponding characterizations and bounds are presented in Theorems~\ref{thm:dist-imagEJR},~\ref{thm:dist_singEJR_i}, and~\ref{thm:dist-hiEJR}.
%More precisely, Theorem~\ref{thm:dist-imagEJR} provides a lower bound for the distance to the imaginary axis, Theorem~\ref{thm:dist_singEJR_i} yields an exact expression for the distance to singularity, and Theorem~\ref{thm:dist-hiEJR} establishes an upper bound for the distance to higher index.
The computed values of these quantities are summarized in the Table~\ref{tab:example1a}.
    %
    \begin{table}[h!]
        \centering
        \renewcommand{\arraystretch}{1.5}
        \begin{tabular}{cccc}
            \hline
            \textbf{$\mathbb  S$}
            & \textbf{l.b. to} $\mathbf{d_{im}^{\mathbb S}(E,J,R)}$
            & $\mathbf{d_{sing}^{\mathbb S}(E,J,R)}$
            & \textbf{u.b. to} $\mathbf{d_{hi}^{\mathbb S}(E,J,R)}$ \\

            & Theorem~\ref{thm:dist-imagEJR}
            & Theorem~\ref{thm:dist_singEJR_i}
            & Theorem~\ref{thm:dist-hiEJR}\\
            \hline

            $\mathcal S_d(E,J,R)$
            & 2.1984
            & 6.3032
            & 8.4024 \\

            $\mathcal S_i(E,J,R)$
            & 1.1467
            & 2.3141
            & 1.1376 \\
            \hline
        \end{tabular}
        \caption{Various structured distances and bounds.}
        \label{tab:example1a}
    \end{table}
    %
    Each of the three distances,
    %i.e., distance to imaginary axis, distance to singularity and the distance to higher index
provides an upper bound for the distance to robust instability. Hence the distance to the nearest system that is not robustly asymptotically stable is bounded from above by $1.1376$ and achieved for a system where the index changes.
%Moreover, the distance to robust instability in this case will be given by the distance to higher index, as the distance to singularity and the lower bound to the distance to imaginary axis are larger than the distance to higher index.

As expected, %it is demonstrated in %Tables~\ref{tab:example1} and~\ref{tab:example1a} that
the structured distances obtained when perturbations are restricted to $J$ and $R$ are consistently larger than those obtained when perturbations in all three matrices $E$, $J$, and $R$ are permitted.
%his behavior is expected, since allowing perturbations in $E$ increases the admissible perturbation set and thus provides additional degrees of freedom.
}
\end{example}

%MehM19_example1.mat
\begin{example}{\rm
    Consider the DC power network example from~\cite{MehM19}. The system can be written as dHDAE
    %
    \[
    E\dot{x} = (J - R)x,
    \]
    where the matrices $E, J, R \in \mathbb{R}^{5\times 5}$ are given by
    \[
    E = \text{diag}(L, C_1, C_2, 0, 0),\quad
    J = \begin{bmatrix}
    0 & -1 & 1 & 0 & 0 \\
    1 & 0 & 0 & -1 & 0 \\
    -1 & 0 & 0 & 0 & -1 \\
    0 & 1 & 0 & 0 & 0 \\
    0 & 0 & 1 & 0 & 0
    \end{bmatrix},\quad
    R = \begin{bmatrix}
    R_L & 0 & 0 & 0 & 0 \\
    0 & 0 & 0 & 0 & 0 \\
    0 & 0 & 0 & 0 & 0 \\
    0 & 0 & 0 & R_G & 0 \\
    0 & 0 & 0 & 0 & R_R
    \end{bmatrix},
    \]
    %
    with inductance $L>0$, capacitances $C_1,C_2>0$, and resistances $R_G,R_L,R_R>0$. Choose the parameters \(L, C_1, C_2, R_L, R_G, R_R\) randomly using the \texttt{randn} command in \textsc{Matlab}. Using the criteria from Section~\ref{sec:distalg},  we have verified that the resulting dHDAE system is robustly asymptotically stable. The generalized eigenvalues of the matrix pencil $(E, J-R)$ are $-37.1617, -2.7775 + 4.0447i, -2.7775 - 4.0447i, \infty, \infty$.  Hence, all finite eigenvalues lie strictly in the open left half-plane, while two infinite eigenvalues arise due to the algebraic constraints. The system is regular and of index one.

    We computed the various structured distances introduced in Section~\ref{sec:partial},  while perturbing only $J$ and $R$. The distances are summarized in Table~\ref{tab:example2}.
    %
    \begin{table}[h!]
        \centering
        \renewcommand{\arraystretch}{1.5}
        \begin{tabular}{cccc}
            \hline
            \textbf{$\mathbb  S$}
            & $\mathbf{d_{im}^{\mathbb S}(J,R)}$
            & $\mathbf{d_{sing}^{\mathbb S}(J,R)}$
            & $\mathbf{d_{hi}^{\mathbb S}(J,R)}$ \\

            & Theorem~\ref{thm:dist-imJR}
            & Theorem~\ref{thm:dist-singJR}
            & Theorem~\ref{thm:dist-hiJR}\\
            \hline
            $\mathcal S_d(J,R)$
            & 1.1043
            & 1.6714
            & 1.3393     \\

            $\mathcal S_i(J,R)$
            & 0.9759
            & 1.6714
            & 1.3393 \\
            \hline
        \end{tabular}
        \caption{Various structured distances.}
        \label{tab:example2}
    \end{table}
%
%Recall that the distance to robust instability is the minimum of the three critical distances, namely the distance to the imaginary axis, the distance to singularity, and the distance to higher index.
Hence, the distance to the boundary of the region of robust asymptotic stability with respect to the perturbation set $\mathcal S_d(J,R)$ is $d_{inst}^{\mathcal S_d}(J,R)=1.1043$, while with respect to $\mathcal S_i(J,R)$ it is $d_{inst}^{\mathcal S_d}(J,R)=0.9759$. \\

The computed  structured distances respectively bounds with respect to perturbations from the full structure-preserving sets $\mathcal S_d(E,J,R) \text{ and }\mathcal S_i(E,J,R)$ %where all three matrices $E$, $J$, and $R$ are allowed to vary. The computed values of these quantities
are summarized in Table~\ref{tab:example2a}.
    %
    \begin{table}[h!]
        \centering
        \renewcommand{\arraystretch}{1.5}
        \begin{tabular}{cccc}
            \hline
            \textbf{$\mathbb  S$}
            & \textbf{l.b. to} $\mathbf{d_{im}^{\mathbb S}(E,J,R)}$
            & $\mathbf{d_{sing}^{\mathbb S}(E,J,R)}$
            & \textbf{u.b. to} $\mathbf{d_{hi}^{\mathbb S}(E,J,R)}$ \\

            & Theorem~\ref{thm:dist-imagEJR}
            & Theorem~\ref{thm:dist_singEJR_i}
            & Theorem~\ref{thm:dist-hiEJR}\\
            \hline

            $\mathcal S_d(E,J,R)$
            & 1.1769
            & 1.1061
            & 2.1160 \\

            $\mathcal S_i(E,J,R)$
            & 0.0192
            & 1.0015
            & 1.0192 \\
            \hline
        \end{tabular}
        \caption{Various structured distances and bounds.}
        \label{tab:example2a}
    \end{table}
    %
Each of the three distances
%, i.e., distance to imaginary axis, distance to singularity and the distance to higher index
provides an upper bound for the distance to nearest system that is not robust asymptotically stable. }
\end{example}

%{\color{blue}
%\begin{enumerate}
%    \item I updated the Example 28 in Numerical section by adding some more reasoning and including the $\mathcal S_i(E,J,R)$ case. The changes are mentioned in blue color. I believe that the paragraph before the blue paragraph in Example 28 can be merged into it. As of now, I have kept them separate.
%    \item I also updated the Example 29 in Numerical section by including the $\mathcal S_i(E,J,R)$ case.
%    \item I will soon create a project file in Gitlab and will upload all the Matlab codes there.
%\end{enumerate}}
%{\color{blue} We can compute various distances while perturbing all three matrices with respect to only $\mathcal S_i(E,J,R)$. As the distance to imaginary axis with respect to $\mathcal S_d(E,J,R)$ is a constrained optimization problem. We don't have any numerical algorithm to tackle this and hence it can me mentioned as a future work.}{\color{red} VM: I suggest to add the results both examples with  $\mathcal S_i(E,J,R)$ and that you look into the use of standard constrained optimization codes from matlab.}

%{\color{red} Should I include an example to compute distance to robust instability with respect to perturbations from the set $\mathcal S_i(E,J,R)$. Yes definitely }
\section{Conclusions}\label{sec:conclusion}

For semidissipative Hamiltonian differential-algebraic systems and the associated matrix pencils, we have presented explicit characterizations as well as bounds for the distance to the nearest system with purely imaginary eigenvalues, the nearest system of index higher than one and the nearest  singular system. Questions of further research are  the equality cases when we only have lower and upper bounds, as well as  appropriate numerical optimization methods to compute these distances. The extension to large-scale systems will require a combination with model reduction techniques.


%%%%%%%%%%%%%%%%%%%%%%%%%%%%%%%%%%%%%%%%%%%%%%%%%%%%%%%%%%%%%%%%%%%%%%%%%%%%%%%%
% *** CrediT statement ***                                                     %
%%%%%%%%%%%%%%%%%%%%%%%%%%%%%%%%%%%%%%%%%%%%%%%%%%%%%%%%%%%%%%%%%%%%%%%%%%%%%%%%

\section*{Author contributions}%
\addcontentsline{toc}{section}{Author contributions}
 Peter Benner performed writing - review and editing, and funding acquisition. Volker Mehrmann and Punit Sharma performed writing - review and editing, and conceptualization.
 Anshul Prajapati performed writing - original draft, conceptualization and numerical implementation.

%%%%%%%%%%%%%%%%%%%%%%%%%%%%%%%%%%%%%%%%%%%%%%%%%%%%%%%%%%%%%%%%%%%%%%%%%%%%%%%%
% *** ACKNOWLEDGEMENTS ***                                                     %
%%%%%%%%%%%%%%%%%%%%%%%%%%%%%%%%%%%%%%%%%%%%%%%%%%%%%%%%%%%%%%%%%%%%%%%%%%%%%%%%

\section*{Acknowledgments}%
\addcontentsline{toc}{section}{Acknowledgments}
%
Anshul Prajapati acknowledges the Max Planck Institute for support through a postdoctoral fellowship. Punit Sharma acknowledges the support of the SERB-CRG grant (CRG/2023/003221) by Government of India.


\begin{thebibliography}{10}

\bibitem{AchAM21}
F.~Achleitner, A.~Arnold, and V.~Mehrmann.
\newblock Hypocoercivity and controllability in linear semi-dissipative {ODEs}
  and {DAEs}.
\newblock {\em ZAMM --- Zeitschrift f{\"u}r Angewandte Mathematik und
  Mechanik}, 103:e202100171, 2021.

\bibitem{AliMM20}
N.~Aliyev, V.~Mehrmann, and E.~Mengi.
\newblock Approximation of stability radii for large-scale dissipative
  {H}amiltonian systems.
\newblock {\em Advances in Computational Mathematics}, 46(1):6, 2020.

\bibitem{AltMU21}
R.~Altmann, V.~Mehrmann, and B.~Unger.
\newblock Port-{H}amiltonian formulations of poroelastic network models.
\newblock {\em Mathematical and Computer Modelling of Dynamical Systems},
  27:429--452, 2021.

\bibitem{BagGS21}
M.~K. Baghel, N.~Gillis, and P.~Sharma.
\newblock Characterization of the dissipative mappings and their application to
  perturbations of dissipative-{H}amiltonian systems.
\newblock {\em Numerical Linear Algebra with Applications}, 28(6):e2402, 2021.

\bibitem{BeaMXZ18}
C.~Beattie, V.~Mehrmann, H.~Xu, and H.~Zwart.
\newblock Linear port-{H}amiltonian descriptor systems.
\newblock {\em Mathematics of Control Signals and Systems}, 30(4):17, 2018.

\bibitem{BenMPS25}
P.~Benner, V.~Mehrmann, A.~Prajapati, and P.~Sharma.
\newblock Computation of structured stability radii for
  dissipative-{H}amiltonian systems, 2025.
\newblock \href {http://arxiv.org/abs/2511.14935} {\path{arXiv:2511.14935}}.

\bibitem{BorKMS14}
S.~Bora, M.~Karow, C.~Mehl, and P.~Sharma.
\newblock Structured eigenvalue backward errors of matrix pencils and
  polynomials with {H}ermitian and related structures.
\newblock {\em SIAM Journal on Matrix Analysis and Applications},
  35(2):453--475, 2014.

\bibitem{Bye88}
R.~Byers.
\newblock A bisection method for measuring the distance of a stable matrix to
  the unstable matrices.
\newblock {\em SIAM Journal on Scientific and Statistical Computing},
  9(5):875--881, 1988.

\bibitem{ByeHM98}
R.~{Byers}, C.~{He}, and V.~{Mehrmann}.
\newblock Where is the nearest non-regular pencil?
\newblock {\em Linear Algebra and its Applications}, 285(1-3):81--105, 1998.

\bibitem{ByeN93}
R.~Byers and N.~K. Nichols.
\newblock On the stability radius of a generalized state-space system.
\newblock {\em Linear Algebra and its Applications}, 188:113--134, 1993.

\bibitem{DalV98}
M.~Dalsmo and A.~van~der Schaft.
\newblock On representations and integrability of mathematical structures in
  energy-conserving physical systems.
\newblock {\em SIAM Journal on Control and Optimization}, 37(1):54--91, 1998.

\bibitem{DuLM13}
N.~H. {Du}, V.~H. {Linh}, and V.~{Mehrmann}.
\newblock Robust stability of differential-algebraic equations.
\newblock In {\em Surveys in Differential-Algebraic Equations {I}}, pages
  63--95. Berlin: Springer, 2013.

\bibitem{EmmM13}
E.~{Emmrich} and V.~{Mehrmann}.
\newblock Operator differential-algebraic equations arising in fluid dynamics.
\newblock {\em Computer Methods in Applied Mathematics}, 13(4):443--470, 2013.

\bibitem{Gan59a}
F.R. Gantmacher.
\newblock {\em Theory of Matrices}, volume~1.
\newblock Chelsea, New York, 1959.

\bibitem{GilMS18}
N.~Gillis, V.~Mehrmann, and P.~Sharma.
\newblock Computing nearest stable matrix pairs.
\newblock {\em Numerical Linear Algebra and its Applications}, 25:e2153, 2018.

\bibitem{GraMQSV16}
N.~Gr{\"a}bner, V.~Mehrmann, S.~Quraishi, C.~Schr{\"o}der, and U.~Von~Wagner.
\newblock Numerical methods for parametric model reduction in the simulation of
  disk brake squeal.
\newblock {\em ZAMM --- Zeitschrift f{\"u}r Angewandte Mathematik und
  Mechanik}, 96(12):1388--1405, 2016.

\bibitem{GudLMS22}
C.~G{\"u}d{\"u}c{\"u}, J.~Liesen, V.~Mehrmann, and D.~Szyld.
\newblock On non-{H}ermitian positive (semi)definite linear algebraic systems
  arising from dissipative {H}amiltonian {DAE}s.
\newblock {\em {SIAM} Journal on Scientific Computing}, 44:A2871--A2894, 2022.

\bibitem{GugL25}
N.~Guglielmi and C.~Lubich.
\newblock Matrix nearness problems and eigenvalue optimization, 2025.
\newblock \href {http://arxiv.org/abs/2503.14750} {\path{arXiv:2503.14750}}.

\bibitem{GugLM17}
N.~Guglielmi, C.~Lubich, and V.~Mehrmann.
\newblock On the nearest singular matrix pencil.
\newblock {\em {SIAM} Journal on Matrix Analysis and Applications},
  38(3):776--806, 2017.

\bibitem{GugM22}
N.~Guglielmi and V.~Mehrmann.
\newblock Computation of the nearest structured matrix triplet with common null
  space.
\newblock {\em Electronic Transactions on Numerical Analysis}, 55:508--531,
  2022.

\bibitem{JacZ12}
B.~Jacob and H.~J. Zwart.
\newblock {\em Linear Port-Hamiltonian Systems on Infinite-dimensional Spaces}.
\newblock Operator Theory: Advances and Applications. Springer, 2012.

\bibitem{LuPSB25}
D.~Lu, A.~Prajapati, P.~Sharma, and S.~Bora.
\newblock Eigenvalue backward errors of {R}osenbrock systems and optimization
  of sums of {R}ayleigh quotients.
\newblock {\em SIAM Journal on Matrix Analysis and Applications},
  46(2):1301--1327, 2025.

\bibitem{MacMT08}
D.~S. Mackey, N.~Mackey, and F.~Tisseur.
\newblock Structured mapping problems for matrices associated with scalar
  products. {P}art {I}: {L}ie and {J}ordan algebras.
\newblock {\em SIAM Journal on Matrix Analysis and Applications},
  29(4):1389--1410, 2008.

\bibitem{Mar86}
N.~Martins.
\newblock Efficient eigenvalue and frequency response methods applied to power
  system small-signal stability studies.
\newblock {\em IEEE Transactions on Power Systems}, 1(1):217--224, 1986.

\bibitem{MarL90}
N.~Martins and L.~T.~G. Lima.
\newblock Determination of suitable locations for power system stabilizers and
  static {VAR} compensators for damping electromechanical oscillations in large
  scale power systems.
\newblock {\em IEEE Transactions on Power Systems}, 5(4):1455--1469, 1990.

\bibitem{MarPR07}
N.~Martins, P.~C. Pellanda, and J.~Rommes.
\newblock Computation of transfer function dominant zeros with applications to
  oscillation damping control of large power systems.
\newblock {\em IEEE Transactions on Power Systems}, 22(4):1657--1664, 2007.

\bibitem{MehMS16}
C.~Mehl, V.~Mehrmann, and P.~Sharma.
\newblock Stability radii for linear {H}amiltonian systems with dissipation
  under structure-preserving perturbations.
\newblock {\em SIAM Journal on Matrix Analysis and Applications},
  37(4):1625--1654, 2016.

\bibitem{MehMS17}
C.~Mehl, V.~Mehrmann, and P.~Sharma.
\newblock Stability radii for real linear {H}amiltonian systems with perturbed
  dissipation.
\newblock {\em BIT Numerical Mathematics}, 57:811--843, 2017.

\bibitem{MehMW18}
C.~{Mehl}, V.~{Mehrmann}, and M.~{Wojtylak}.
\newblock Linear algebra properties of dissipative {H}amiltonian descriptor
  systems.
\newblock {\em SIAM Journal on Matrix Analysis and Applications},
  39:1489--1519, 2018.

\bibitem{MehMW21}
C.~Mehl, V.~Mehrmann, and M.~Wojtylak.
\newblock Distance problems for dissipative {H}amiltonian systems and related
  matrix polynomials.
\newblock {\em Linear Algebra and its Applications}, 623:335--366, 2021.

\bibitem{MehM19}
V.~Mehrmann and R.~Morandin.
\newblock Structure-preserving discretization for port-{Hamiltonian} descriptor
  systems.
\newblock In {\em 2019 IEEE 58th Conference on Decision and Control (CDC)},
  pages 6863--6868. IEEE, 2019.

\bibitem{MehU23}
V.~Mehrmann and B.~Unger.
\newblock Control of port-{H}amiltonian differential-algebraic systems and
  applications.
\newblock {\em Acta Numerica}, pages 395--515, 2023.

\bibitem{MehX15}
V.~Mehrmann and H.~Xu.
\newblock Structure preserving deflation of infinite eigenvalues in structured
  pencils.
\newblock {\em Electronic Transactions on Numerical Analysis}, 44:1--24, 2015.

\bibitem{OrtVME02}
R.~Ortega, A.~van~der Schaft, B.~Maschke, and G.~Escobar.
\newblock Interconnection and damping assignment passivity-based control of
  port-controlled {H}amiltonian systems.
\newblock {\em Automatica}, 38(4):585--596, 2002.

\bibitem{PraS22}
A.~Prajapati and P.~Sharma.
\newblock Estimation of structured distances to singularity for matrix pencils
  with symmetry structures: A linear algebra--based approach.
\newblock {\em SIAM Journal on Matrix Analysis and Applications},
  43(2):740--763, 2022.

\bibitem{RomM09}
J.~Rommes and N.~Martins.
\newblock Exploiting structure in large-scale electrical circuit and power
  system problems.
\newblock {\em Linear Algebra and its Applications}, 431(3-4):318--333, 2009.

\bibitem{Sch90}
W.~Schiehlen.
\newblock {\em Multibody Systems Handbook}.
\newblock Springer, Heidelberg, Germany, 1990.

\bibitem{Ves11}
K.~Veseli{\'c}.
\newblock {\em Damped Oscillations of Linear Systems: a Mathematical
  Introduction}.
\newblock Springer Science \& Business Media, 2011.

\end{thebibliography}

%\addcontentsline{toc}{section}{References}
%\bibliographystyle{plainurl}
%\bibliography{exampleref}

\end{document}

\appendix
\section{Appendix A}


{\bf Staircase Algorithm for triple~$(E,J,R)$ with $E=E^*\geq 0$ $J=-J^*$ and $R=R^*\geq 0$.}

{\em Step 1}\\
i) Perform a spectral decomposition of~$E$ such that
\[
E =U_E \begin{bmatrix} E_{1,1} & 0 \\ 0 & 0 \end{bmatrix} U_E^*,
\]
with $U_E\in\mathbb C^{n,n}$ unitary, $E_{1,1}\in \mathbb C^{\tilde n_1,  \tilde n_1}$ positive definite or $\tilde n_1=0$.\\
%
ii) Set $P := U_E$,
$
J:= U^*_E\ J\ U_E
=\begin{bmatrix}
 J_{1,1} & -J_{2,1}^* \\
 J_{2,1} & J_{2,2}
 \end{bmatrix}$, $R:= U^*_E\ R\ U_E
=\begin{bmatrix}
 R_{1,1} & R_{2,1}^* \\
 R_{2,1} & R_{2,2} \end{bmatrix}$, and $
E:=U^*_E\ E\ U_E$.
%
{\em Step 2}\\
IF {$\tilde n_1<n$}\\
iii)  Compute full rank decomposition $J_{2,2}-R_{2,2}\in\mathbb C^{(n-\tilde n_1)\times (n-\tilde n_1)}$ such that
\[
P_{2,2}^*\ (J_{2,2}-R_{2,2})\ P_{2,2}
=\begin{bmatrix} \Sigma_{2,2} & 0 \\ 0 & 0\end{bmatrix},
\]
with $\Sigma_{2,2}\in \mathbb C^{\tilde n_2,\tilde n_2}$ invertible or $\tilde n_2=0$.\\
END IF
\\
iv)  Set
\[
P_2:=\begin{bmatrix} I& 0 \\ 0 & P_{2,2} \end{bmatrix} \in\mathbb C^{n,n},\qquad P:=P  P_2.
\]
v) Set $E:=P_2^*\ E\ P_2$,
\begin{eqnarray*}
J &:=& P_2^*\ J\ P_2
=: \left[ \begin{array}{c|cc}
 J_{1,1} & -J_{2,1}^* & -J_{3,1}^*  \\
 \hline %\\[-3.6mm]
 J_{2,1} & J_{2,2} & 0 \\
 J_{3,1}& 0 & 0
\end{array}\right],\\
R &:=& P_2^*\ R\ P_2
=: \left[ \begin{array}{c|cc}
 R_{1,1} & R_{2,1}^* & 0 \\
 \hline %\\[-3.6mm]
 R_{2,1} & R_{2,2} & 0 \\
 0 & 0 & 0
\end{array}\right],
\end{eqnarray*}
with $J_{2,2}-R_{2,2}=\Sigma_{2,2}$.
(The lines indicate the partitioning of the block matrices~$J$ and~$R$ in the previous step.
Note that the positive semi-definiteness of the Hermitian matrix~$R$ implies the~$0$ structure in~$R$.) \\
%
{\em Step 3}\\
vi) Define $\tilde n_3 :=n -\tilde n_1 -\tilde n_2$.\\
IF {$\tilde n_3>0$}\\
vii) Perform an SVD of $J_{3,1}$ such that
\[
J_{3,1} = U_{3,1}^* \begin{bmatrix} \Sigma_{3,1} & 0\\ 0 & 0 \end{bmatrix} V_{3,1} \in \mathbb C^{\tilde n_3 \times \tilde n_1}\,,
\]
with~$\Sigma_{3,1}\in\mathbb R^{n_1 \times n_1}$ nonsingular diagonal or $n_1 =0$.\\
END IF\\
viii) Set
\[
P_3 :=\begin{bmatrix}
 V_{3,1} & & \\
  & I & \\
  & & U_{3,1}
\end{bmatrix} \in\mathbb C^{n,n},\qquad
P:= P P_3.
\]
ix) Set $\check E
:= P_3^*\ E\ P_3$, $\check J
:= P_3^*\ J\ P_3$, $\check R
:= P_3^*\ R\ P_3$ such that
\[
\check J
%:= P_3^*\ J\ P_3
=:\left[ \begin{array}{cc|c|cc}
J_{1,1}  &  -J_{2,1}^*  &  -J_{3,1}^*  &  -J_{4,1}^*  &  0 \\
J_{2,1}  &  J_{2,2}  &  -J_{3,2}^*  &  0  &  0 \\
\hline
J_{3,1}  &  J_{3,2}  &  J_{3,3}  &  0  &  0 \\
\hline
J_{4,1}  &  0  &  0  &  0  &  0 \\
0  &  0  &  0  &  0  &  0
\end{array}\right],\
%
\check R
%:= P_3\ R\ P_3^*
=:\left[ \begin{array}{cc|c|cc}
R_{1,1}  &  R_{2,1}^*  &  R_{3,1}^* &  0  &  0 \\
R_{2,1}  &  R_{2,2}  &  R_{3,2}^*  &  0  &  0 \\
\hline
R_{3,1}  &  R_{3,2}  &  R_{3,3}    &  0  &  0 \\
\hline
0  &  0  &  0  &  0  &  0 \\
0  &  0  &  0  &  0  &  0
\end{array}\right],
\]
%
\[
\check E
%:= P_3\ E\ P_3^*
=: \left[ \begin{array}{cc|c|cc}
 E_{1,1} &    E_{2,1}^*    &  0 &0 & 0 \\
 E_{2,1} & E_{2,2} &  0    &   0      & 0 \\
\hline
 0  & 0  & 0  & 0  & 0 \\
\hline
 0  & 0  & 0  & 0  & 0 \\
 0  & 0  & 0  & 0  & 0
\end{array}\right],
\]
%
which are of the desired form with $n_2:=\tilde n_1 -n_1$, $n_3 :=\tilde n_2$, $n_4:=n_1$, $n_5:=\tilde n_3 -n_4$.


%Towards the end, there can be an unnumbered section for the acknowledgments,
%stating people and organizations who influenced or funded the work in the paper.
%Note that this section is just created via
%\begin{verbatim}
%\section*{Acknowledgments}%
%\addcontentsline{toc}{section}{Acknowledgments}
%\end{verbatim}


%%%%%%%%%%%%%%%%%%%%%%%%%%%%%%%%%%%%%%%%%%%%%%%%%%%%%%%%%%%%%%%%%%%%%%%%%%%%%%%%
% *** REFERENCES ***                                                           %
%%%%%%%%%%%%%%%%%%%%%%%%%%%%%%%%%%%%%%%%%%%%%%%%%%%%%%%%%%%%%%%%%%%%%%%%%%%%%%%%
\subsection{Stability of general DAEs}\label{sec:stability}
In contrast to ordinary differential equations (ODE)  the solvability theory of DAEs is more complex because more smoothness may be required and one has to require consistency of the initial condition to have a classical (continuously differentiable) solution.
An initial condition  is called {\em consistent}
if the corresponding initial value problem has at least one classical solution. See \cite{KunM24} for a detailed analysis of DAEs.

To analyze the (asymptotic) stability of DAEs we recall the following classical stability concepts, see \cite{HinP05} for ordinary
differential equations
%
\begin{equation}\label{ODE}
\dot x(t) =f(t,x(t)),\quad t\in \mathbb I,
\end{equation}
%
where $\mathbb I$ is either a finite interval $[t_0,t_1]$ or infinite interval $t\geq t_0$.
\begin{definition}\label{stabtraj}
A solution~$x:t\mapsto x(t;t_0,x_0)$ of (\ref{ODE}) with initial condition $x(t_0)=x_0$ is called
%
\begin{enumerate}
\item
\emph{stable} if for every $\varepsilon >0$ there exists $\delta>0$ such that
\begin{enumerate}
\item
the initial value problem (\ref{ODE}) with initial condition
$x(t_0)=\hat x_0$ is solvable on~${\mathbb I}$ for all $\hat x_0\in {\mathbb K}^n$
with $\| \hat x_0 -x_0 \|<\delta$;
\item
the solution $x(t;t_0,\hat x_0)$ satisfies
$\| x(t;t_0,\hat x_0) -x(t;t_0,x_0) \|<\varepsilon$ on~${\mathbb I}$.
\end{enumerate}
\item
\emph{asymptotically stable} if
it is stable and there exists $\varrho>0$ such that
\begin{enumerate}
\item
the initial value problem (\ref{ODE}) with initial condition
$x(t_0)=\hat x_0$ is solvable on~${\mathbb I}$ for all $\hat x_0\in {\mathbb K}^n$
with $\| \hat x_0 -x_0 \|<\varrho$;
\item
the solution $x(t;t_0,\hat x_0)$ satisfies
$\lim_{t\to \infty}\| x(t;t_0,\hat x_0) -x(t;t_0,x_0) \|=0$.
\end{enumerate}
\end{enumerate}
\end{definition}
%
Note that by a shift one can always transform the ODE (\ref{ODE}) in such a way
that a given solution $x(t;t_0,x_0)$ is mapped to the trivial solution. We therefore assume always that the equilibrium solution is the trivial solution, and also that $t_0=0$.


One can immediately extend Definition~\ref{stabtraj} verbatim
to regular DAEs. However, one has to make sure that the perturbed initial condition is still consistent. For a detailed analysis of the stability of general DAEs, see \cite{KunM24,MehU23}.

